% ============================================================
%  Rapport PFA 2025-2026 * SecureWatch
%  Compilateur : pdflatex (deux passages)
%  Encodage   : UTF-8
% ============================================================
\documentclass[12pt, a4paper]{report}

% ── Encodage & langue ────────────────────────────────────────
\usepackage[utf8]{inputenc}
\usepackage[T1]{fontenc}
\usepackage[french]{babel}

% ── Mise en page ─────────────────────────────────────────────
\usepackage[
  top=2.5cm, bottom=2.5cm,
  left=2.8cm, right=2.5cm,
  headheight=15pt
]{geometry}
\usepackage{setspace}
\onehalfspacing

% ── Graphiques & couleurs ─────────────────────────────────────
\usepackage{graphicx}
\usepackage[dvipsnames,table]{xcolor}
\usepackage{float}
\usepackage{caption}
\usepackage{subcaption}
\usepackage{tikz}
\graphicspath{{./}{screenshots/}}

% Couleurs charte EMSI
\definecolor{emsigreen}{RGB}{0, 130, 60}
\definecolor{emsiref}{RGB}{180, 0, 0}

% ── Typographie ───────────────────────────────────────────────
\usepackage{lmodern}
\usepackage{microtype}
\usepackage{parskip}

% ── En-têtes et pieds de page ────────────────────────────────
\usepackage{fancyhdr}
\setlength{\headheight}{14.5pt}
\pagestyle{fancy}
\fancyhf{}
\fancyhead[L]{\small\textit{SecureWatch / PFA 2025-2026}}
\fancyhead[R]{\small\textit{\nouppercase{\leftmark}}}
\fancyfoot[C]{\thepage}
\renewcommand{\headrulewidth}{0.4pt}

% Supprime "Chapitre N." du mark * garde uniquement le titre du chapitre
\renewcommand{\chaptermark}[1]{\markboth{#1}{}}
% Supprime la numérotation de section dans le mark droit
\renewcommand{\sectionmark}[1]{\markright{#1}}

% Style plain pour les premières pages de chapitre (évite le numéro en haut)
\fancypagestyle{plain}{%
  \fancyhf{}%
  \fancyfoot[C]{\thepage}%
  \renewcommand{\headrulewidth}{0pt}%
}

% ── Liens hypertexte ─────────────────────────────────────────
\usepackage[
  colorlinks=true,
  linkcolor=NavyBlue,
  urlcolor=NavyBlue,
  citecolor=NavyBlue
]{hyperref}

% ── Code source ──────────────────────────────────────────────
\usepackage{listings}
\usepackage{lstautogobble}

\definecolor{codebg}{RGB}{248,249,250}
\definecolor{codeframe}{RGB}{222,226,230}
\definecolor{keyword}{RGB}{0,119,170}
\definecolor{string}{RGB}{102,153,0}
\definecolor{comment}{RGB}{153,153,153}

\lstdefinestyle{securewatch}{
  backgroundcolor=\color{codebg},
  frame=single,
  framerule=0.4pt,
  rulecolor=\color{codeframe},
  basicstyle=\ttfamily\footnotesize,
  keywordstyle=\color{keyword}\bfseries,
  stringstyle=\color{string},
  commentstyle=\color{comment}\itshape,
  breaklines=true,
  breakatwhitespace=true,
  showstringspaces=false,
  tabsize=2,
  autogobble=true,
  numbers=left,
  numberstyle=\tiny\color{gray},
  numbersep=8pt,
  xleftmargin=12pt,
  captionpos=b
}

\lstdefinestyle{bash}{
  style=securewatch,
  language=bash,
  morekeywords={docker-compose,python,curl,pip}
}

\lstdefinestyle{yaml}{
  style=securewatch,
  language=,
  morestring=[b]{"},
  morestring=[b]{'},
  morecomment=[l]{\#},
  morekeywords={true,false,null}
}

\lstset{style=securewatch}

% ── Tableaux ─────────────────────────────────────────────────
\usepackage{booktabs}
\usepackage{tabularx}
\usepackage{longtable}
\usepackage{multirow}
\usepackage{array}

% ── Boîtes d'information ─────────────────────────────────────
\usepackage{tcolorbox}
\tcbuselibrary{skins, breakable}

\newtcolorbox{infobox}[1][]{
  enhanced, breakable,
  colback=blue!5, colframe=NavyBlue!60,
  fonttitle=\bfseries, title=#1,
  left=6pt, right=6pt, top=4pt, bottom=4pt
}

\newtcolorbox{alertbox}[1][]{
  enhanced, breakable,
  colback=red!5, colframe=red!60,
  fonttitle=\bfseries, title=#1,
  left=6pt, right=6pt, top=4pt, bottom=4pt
}

\newtcolorbox{successbox}[1][]{
  enhanced, breakable,
  colback=green!5, colframe=green!50!black!60,
  fonttitle=\bfseries, title=#1,
  left=6pt, right=6pt, top=4pt, bottom=4pt
}

% ── Divers ───────────────────────────────────────────────────
\usepackage{enumitem}
\usepackage{titlesec}
\usepackage{amsmath}
\usepackage{amssymb}
\usepackage{pifont}

\newcommand{\cmark}{\ding{51}}
\newcommand{\xmark}{\ding{55}}

% ── Placeholder automatique pour les screenshots ─────────────
% Cherche dans cet ordre :
%   1. screenshots/fichier.png   2. screenshots/fichier.jpg
%   3. fichier.png (meme dossier que le .tex)
%   4. fichier.jpg (meme dossier que le .tex)
% Si aucun fichier trouve -> boite grise placeholder
\newcommand{\screenshot}[4][0.95\textwidth]{%
  \begin{figure}[H]
    \centering
    \IfFileExists{screenshots/#2.png}{%
      \includegraphics[width=#1]{screenshots/#2.png}%
    }{%
    \IfFileExists{screenshots/#2.jpg}{%
      \includegraphics[width=#1]{screenshots/#2.jpg}%
    }{%
    \IfFileExists{#2.png}{%
      \includegraphics[width=#1]{#2.png}%
    }{%
    \IfFileExists{#2.jpg}{%
      \includegraphics[width=#1]{#2.jpg}%
    }{%
      \fcolorbox{gray!60}{gray!10}{%
        \parbox[c][5cm][c]{\dimexpr#1 - 2\fboxsep - 2\fboxrule\relax}{%
          \centering
          {\color{gray!50}\fontsize{36}{36}\selectfont\ding{94}}\\[0.4cm]
          {\ttfamily\small\color{gray!80} screenshots/#2.png}\\[0.3cm]
          {\small\itshape\color{gray} Screenshot a inserer}%
        }%
      }%
    }}}}%
    \caption{#3}
    \label{fig:#4}
  \end{figure}%
}

% ── Style des titres ─────────────────────────────────────────
\titleformat{\chapter}[display]
  {\normalfont\Large\bfseries\color{NavyBlue}}
  {\chaptertitlename\ \thechapter}{10pt}{\Huge}
\titlespacing*{\chapter}{0pt}{-20pt}{20pt}

\titleformat{\section}
  {\normalfont\large\bfseries\color{NavyBlue!80}}
  {\thesection}{1em}{}
\titleformat{\subsection}
  {\normalfont\normalsize\bfseries}
  {\thesubsection}{1em}{}

% ============================================================
%  DOCUMENT
% ============================================================
\begin{document}

% ── Page de garde ────────────────────────────────────────────
\begin{titlepage}
  \thispagestyle{empty}

  \begin{tikzpicture}[remember picture, overlay]
    \fill[emsigreen]
      (current page.north west)
      rectangle
      ([xshift=1.4cm] current page.south west);
    \fill[emsiref]
      (current page.south west)
      rectangle
      ([xshift=1.4cm, yshift=1.4cm] current page.south west);
  \end{tikzpicture}

  \noindent
  \begin{minipage}[c]{0.48\textwidth}
    \raggedright
    \includegraphics[height=1.2cm]{logo_emsi.png}
  \end{minipage}
  \hfill
  \begin{minipage}[c]{0.40\textwidth}
    \raggedleft
    \includegraphics[height=2.2cm]{logo_honoris.png}
  \end{minipage}

  \vspace{1.8cm}

  \begin{center}
    {\large\bfseries Rapport de projet fin d'année}\\[0.4cm]
    {\large\bfseries 4\textsuperscript{ème} année}\\[0.3cm]
    {\Large\bfseries Ingénierie Informatique et Réseaux}\\[0.25cm]
    {\normalsize\bfseries Cybersécurité \& Infrastructures Réseaux}\\[0.6cm]
    {\normalsize Sous le thème}\\[0.3cm]
    \noindent\rule{0.55\textwidth}{0.8pt}
  \end{center}

  \vspace{0.9cm}

  \begin{center}
    {\fontsize{28}{34}\selectfont\bfseries SECUREWATCH}\\[0.5cm]
    {\large\bfseries
      Sécurisation et Surveillance d'une Architecture\\[0.2cm]
      Microservices - Centralisation des Logs \&\\[0.2cm]
      Détection d'Incidents de Sécurité
    }
  \end{center}

  \vspace{0.5cm}

  \begin{center}
    \noindent\rule{0.55\textwidth}{0.8pt}
  \end{center}

  \vspace{1.6cm}

  \noindent
  \begin{minipage}[t]{0.55\textwidth}
    \raggedright
    {\bfseries Réalisé par :}\\[0.3cm]
    EL HAMA Yahya\\[0.15cm]
    AGRAD Hicham\\[0.8cm]
    {\bfseries Encadré par :}\\[0.3cm]
    Mr. Imam ALIHAMIDI
  \end{minipage}

  \vfill

  \begin{center}
    {\normalsize\bfseries ANNEE UNIVERSITAIRE: 2025-2026}
  \end{center}

\end{titlepage}

% ── Dédicaces ─────────────────────────────────────────────────
\chapter*{Dédicaces}
\thispagestyle{empty}
\addcontentsline{toc}{chapter}{Dédicaces}
\markboth{Dédicaces}{Dédicaces}

\vspace*{2cm}

\begin{center}
  \itshape\large

  À nos familles,\\[0.5cm]
  pour leur soutien indéfectible et leurs encouragements\\
  tout au long de notre parcours académique.\\[1.5cm]

  À nos parents,\\[0.5cm]
  qui ont consenti à tant de sacrifices pour nous offrir\\
  la meilleure éducation possible.\\[1.5cm]

  À tous nos amis et camarades de promotion,\\[0.5cm]
  pour les moments partagés, les nuits de travail\\
  et l'entraide qui ont marqué ces années d'études.\\[1.5cm]

  Nous dédions ce travail.
\end{center}

\clearpage

% ── Remerciements ─────────────────────────────────────────────
\chapter*{Remerciements}
\thispagestyle{empty}
\addcontentsline{toc}{chapter}{Remerciements}
\markboth{Remerciements}{Remerciements}

Ce projet n'aurait pas été possible sans l'aide de nombreuses personnes, et nous tenons à les remercier sincèrement.

Notre première pensée va à notre encadrant, \textbf{Monsieur Imam ALIHAMIDI}. Sa disponibilité, ses conseils avisés et sa connaissance approfondie de la sécurité informatique et des architectures distribuées nous ont été précieux à chaque étape. Les orientations qu'il nous a données ont évité bien des erreurs d'aiguillage.

Nous remercions également l'ensemble du \textbf{corps professoral de l'EMSI} pour quatre années d'enseignements de qualité en Ingénierie Informatique et Réseaux. Ce projet est, en grande partie, le fruit de ce qui a été appris en cours.

Nos remerciements vont aussi aux \textbf{membres du jury}, qui ont accepté de consacrer du temps à l'évaluation de ce travail. Leurs retours et recommandations sont une partie essentielle du processus.

Enfin, merci à nos \textbf{familles}. Les nuits longues et les week-ends absorbés par le projet n'ont pas toujours été faciles à vivre pour eux. Leur soutien a compté.

\clearpage

% ── Résumé / Abstract ─────────────────────────────────────────
\chapter*{Résumé}
\addcontentsline{toc}{chapter}{Résumé / Abstract}
\markboth{Résumé}{Résumé}

\section*{Résumé}

Dans un contexte marqué par la multiplication des architectures distribuées et l'augmentation des cybermenaces, la sécurisation des systèmes d'information représente un enjeu critique. Ce projet présente \textbf{SecureWatch}, une plateforme de centralisation des journaux applicatifs et de détection automatique d'incidents de sécurité, déployée localement via Docker Compose.

L'architecture repose sur 14 services containerisés répartis sur trois réseaux Docker isolés : une API Gateway Nginx assurant la terminaison TLS et le rate limiting, deux microservices Flask instrumentés avec Prometheus, une stack ELK complète (Elasticsearch, Logstash, Kibana, Filebeat) pour la centralisation et l'analyse des logs, un cache Redis pour la persistance des compteurs de brute force, un système de monitoring en temps réel (Prometheus, Grafana, Alertmanager) et un tableau de bord interactif développé en JavaScript pur. La plateforme intègre un contrôle d'accès basé sur les rôles (\textbf{RBAC} à trois niveaux : admin, operator, user) avec application côté serveur et restrictions de navigation côté client, un système d'alertes en temps réel par \textbf{Server-Sent Events (SSE)}, une politique de rétention automatique des index Elasticsearch (\textbf{ILM}, 14 jours), ainsi que la détection de neuf types d'incidents incluant les injections SQL et la détection de User-Agents offensifs. Plus de 4\,400 événements sont indexés lors d'une session de simulation complète, avec des alertes déclenchées en moins de 30 secondes après détection.

\medskip
\noindent\textbf{Mots-clés :} Microservices, Docker, ELK Stack, API Gateway, Cybersécurité, Prometheus, Détection d'incidents, Rate Limiting, TLS, Journalisation.

\vspace{1.2cm}
\newpage
\section*{Abstract}

In a context marked by the proliferation of distributed architectures and the increase of cyber threats, securing information systems represents a critical challenge. This project presents \textbf{SecureWatch}, a platform for centralizing application logs and automatically detecting security incidents, deployed locally via Docker Compose.

The architecture relies on 14 containerized services distributed across three isolated Docker networks: an Nginx API Gateway handling TLS termination and rate limiting, two Prometheus-instrumented Flask microservices, a complete ELK stack (Elasticsearch, Logstash, Kibana, Filebeat) for log centralization and analysis, a Redis cache for brute-force counter persistence, a real-time monitoring system (Prometheus, Grafana, Alertmanager), and an interactive dashboard developed in pure JavaScript. The platform incorporates a three-tier \textbf{RBAC} system (admin/operator/user) with server-side enforcement, \textbf{Server-Sent Events (SSE)} for real-time alert delivery, an Elasticsearch \textbf{ILM} policy for automatic 14-day index retention, and detection of nine incident types including SQL injection and offensive User-Agent detection. Results show reliable collection and indexing of over 4,400 security events, with alerts triggered within 30 seconds of attack detection.

\medskip
\noindent\textbf{Keywords:} Microservices, Docker, ELK Stack, API Gateway, Cybersecurity, Prometheus, Incident Detection, Rate Limiting, TLS, Logging.

\clearpage

% ── Table des matières ────────────────────────────────────────
\tableofcontents
\clearpage

% ── Liste des figures ─────────────────────────────────────────
\listoffigures
\clearpage

% ── Liste des tableaux ────────────────────────────────────────
\listoftables
\clearpage

% ── Glossaire / Liste des Abréviations ───────────────────────
\chapter*{Glossaire et Liste des Abréviations}
\addcontentsline{toc}{chapter}{Glossaire et Liste des Abréviations}
\markboth{Glossaire}{Glossaire}

\begin{description}[leftmargin=3.5cm, style=nextline]

  \item[API] \textit{Application Programming Interface} -> Interface de programmation permettant à des applications de communiquer entre elles via des contrats d'échange bien définis.

  \item[API Gateway] Composant logiciel servant de point d'entrée unique vers un ensemble de microservices. Assure le routage, la sécurité (TLS, authentification), le rate limiting et la journalisation.

  \item[CERT] \textit{Computer Emergency Response Team} -> Équipe de réponse aux incidents de sécurité informatique.

  \item[CI/CD] \textit{Continuous Integration / Continuous Deployment} -> Pratiques d'automatisation de l'intégration et du déploiement logiciel.

  \item[CORS] \textit{Cross-Origin Resource Sharing} -> Mécanisme de sécurité du navigateur limitant les requêtes HTTP cross-origin. Nécessite des en-têtes spécifiques côté serveur.

  \item[DDoS] \textit{Distributed Denial of Service} -> Attaque par déni de service distribué visant à rendre un service indisponible.

  \item[Docker] Plateforme de conteneurisation permettant d'encapsuler des applications et leurs dépendances dans des conteneurs isolés et portables.

  \item[ELK Stack] Suite logicielle composée d'Elasticsearch, Logstash et Kibana, utilisée pour la collecte, le traitement et la visualisation de logs à grande échelle.

  \item[ETL] \textit{Extract, Transform, Load} -> Processus d'extraction, transformation et chargement de données d'une source vers une destination.

  \item[Filebeat] Agent léger de collecte de fichiers de logs, composant de la suite Elastic. Surveille les fichiers et transmet leur contenu à Logstash ou Elasticsearch.

  \item[Grafana] Plateforme open-source de visualisation et d'analyse de métriques temporelles. Supporte de nombreuses sources de données dont Prometheus.

  \item[HTTPS] \textit{HyperText Transfer Protocol Secure} -> Extension sécurisée du protocole HTTP utilisant TLS pour chiffrer les communications.

  \item[JWT] \textit{JSON Web Token} -> Standard ouvert (RFC 7519) définissant un format compact et auto-portant pour transmettre des informations de manière sécurisée entre parties.

  \item[Kibana] Interface web de visualisation des données Elasticsearch. Permet l'exploration des logs, la création de dashboards et la recherche via KQL.

  \item[KQL] \textit{Kibana Query Language} -> Langage de requête propriétaire Kibana pour filtrer et rechercher dans les données Elasticsearch.

  \item[Logstash] Pipeline de traitement de données côté serveur, composant ELK. Ingère des données de multiples sources, les transforme et les envoie vers des destinations configurées.

  \item[Microservice] Style architectural décomposant une application en services petits, indépendants et déployables séparément, chacun responsable d'une fonctionnalité métier précise.

  \item[OAuth2] \textit{Open Authorization 2.0} -> Standard ouvert d'autorisation permettant à une application d'accéder à des ressources protégées au nom d'un utilisateur.

  \item[PromQL] \textit{Prometheus Query Language} -> Langage fonctionnel de requête pour Prometheus permettant de sélectionner et agréger des données de séries temporelles.

  \item[Prometheus] Système de monitoring open-source spécialisé dans la collecte de métriques par modèle pull. Standard de facto pour les environnements cloud-native.

  \item[RBAC] \textit{Role-Based Access Control} -> Modèle de contrôle d'accès dans lequel les permissions sont associées à des rôles attribués aux utilisateurs.

  \item[Rate Limiting] Mécanisme de contrôle du débit des requêtes, limitant le nombre d'appels qu'un client peut effectuer par unité de temps.

  \item[REST] \textit{Representational State Transfer} -> Style architectural pour les APIs web, reposant sur les méthodes HTTP (GET, POST, PUT, DELETE).

  \item[SIEM] \textit{Security Information and Event Management} -> Système agrégeant et analysant les données de sécurité de multiples sources pour la détection des menaces en temps réel.

  \item[TLS] \textit{Transport Layer Security} -> Protocole cryptographique assurant la confidentialité et l'intégrité des communications sur un réseau.

  \item[TSDB] \textit{Time Series Database} -> Base de données optimisée pour le stockage et la requête de données horodatées.

  \item[YAML] \textit{YAML Ain't Markup Language} -> Format de sérialisation de données lisible par l'humain, utilisé pour les fichiers de configuration.

\end{description}

\clearpage

% ============================================================
%  INTRODUCTION GÉNÉRALE
% ============================================================
\chapter*{Introduction générale}
\addcontentsline{toc}{chapter}{Introduction générale}
\markboth{Introduction générale}{Introduction générale}

Il y a quelques années encore, déployer une application signifiait généralement lancer un seul serveur, un seul processus. Aujourd'hui, la même application peut mobiliser des dizaines de microservices, chacun tournant dans son propre conteneur, chacun générant ses propres logs, chacun représentant une cible potentielle. Cette évolution a profondément changé la donne en matière de sécurité.

Le problème, c'est que la plupart des équipes n'ont pas les outils pour suivre. Quand un attaquant sonde l'architecture, les traces de son passage sont éparpillées entre les journaux de la gateway, ceux du service d'authentification, ceux de l'API. Reconstituer la séquence des événements prend du temps — parfois trop. Les dommages sont souvent causés bien avant que quiconque ait remarqué quelque chose d'anormal.

\textbf{SecureWatch} est né de ce constat. Réalisé dans le cadre du Projet de Fin d'Année 2025-2026 de la filière Ingénierie Informatique et Réseaux, spécialité Cybersécurité \& Infrastructures Réseaux, ce projet propose une plateforme open-source capable de centraliser les journaux, de détecter les incidents automatiquement et d'alerter en temps réel — le tout déployable en une seule commande sur un poste local.

Ce rapport se divise en quatre chapitres. Le premier pose les bases : contexte, problème et objectifs. Le deuxième explore les technologies disponibles et justifie les choix retenus. Le troisième entre dans le détail de l'implémentation. Le quatrième analyse ce que les simulations d'attaques ont révélé sur l'efficacité du système. Une conclusion tire le bilan et trace les pistes d'amélioration.

% ============================================================
%  CHAPITRE 1 * CADRAGE DU PROJET
% ============================================================
\chapter{Cadrage du Projet}

\section{Introduction}

Avant d'entrer dans les détails techniques, ce chapitre pose les fondations du projet. Pourquoi SecureWatch existe-t-il ? Quel problème concret cherche-t-il à résoudre ? Et comment avons-nous organisé le travail pour y répondre efficacement ?

\section{Contexte}

Les architectures microservices ont changé la manière dont on construit les logiciels. Là où une application monolithique formait un bloc unique, on travaille désormais avec des dizaines de services indépendants, chacun conteneurisé, chacun déployable séparément. Docker et Docker Compose ont démocratisé cette approche, même pour des équipes de petite taille.

Mais cette flexibilité a un coût. Plus il y a de services, plus la surface d'attaque s'étend. Un conteneur mal configuré, une API exposée sans protection, un secret hardcodé dans une image — autant de failles qui peuvent passer inaperçues pendant des semaines. Et dans un environnement distribué, les traces d'une intrusion sont rarement regroupées au même endroit : elles sont dispersées dans les logs de chaque service, rendant la corrélation manuelle pratiquement impossible.

Les chiffres parlent d'eux-mêmes : selon le rapport Verizon DBIR 2024, plus de 40\% des incidents de sécurité impliquent désormais des APIs ou des architectures cloud. Surveiller efficacement ces environnements n'est plus un luxe — c'est une nécessité.

\begin{infobox}[Contexte pédagogique]
Ce projet s'inscrit dans le cadre de la formation en Cybersécurité \& Infrastructures Réseaux de l'EMSI. Il vise à mettre en pratique les concepts théoriques de sécurisation des architectures distribuées, de gestion des accès et de surveillance des systèmes, en utilisant exclusivement des outils open-source.
\end{infobox}

\section{Problématique}

Dans un environnement microservices, les approches de sécurité traditionnelles montrent rapidement leurs limites. Les défis qui se posent concrètement sont les suivants :

\begin{itemize}[leftmargin=*]
  \item \textbf{Logs dispersés} : chaque service écrit ses journaux de manière indépendante dans des fichiers distincts, rendant la corrélation manuelle entre événements pratiquement impossible à l'échelle.
  \item \textbf{Détection réactive} : sans mécanisme de corrélation automatique, les incidents sont souvent détectés après coup, lorsque les dommages sont déjà causés.
  \item \textbf{Absence de protection à l'entrée} : les microservices exposés directement sur le réseau sont vulnérables aux attaques par force brute, aux scans de ressources sensibles et aux abus d'API.
  \item \textbf{Absence de monitoring} : sans métriques temps réel, l'état de santé du système et les anomalies comportementales ne sont pas observables.
  \item \textbf{Fragmentation des outils} : l'absence d'une vue unifiée oblige les équipes de sécurité à consulter de multiples interfaces, augmentant le temps de réponse aux incidents.
\end{itemize}

\textbf{Question centrale} : Comment sécuriser efficacement une architecture microservices locale avec des outils open-source, tout en centralisant la détection d'incidents et en offrant une observabilité complète en temps réel ?

\section{Objectifs}

\begin{successbox}[Objectifs du projet SecureWatch]
\begin{enumerate}[leftmargin=*]
  \item \textbf{Sécuriser les accès} : déployer une API Gateway avec terminaison TLS, rate limiting par zone et isolation réseau pour protéger les microservices de toute exposition directe.
  \item \textbf{Contrôle d'accès par rôles (RBAC)} : implémenter un système à trois niveaux (admin, operator, user) avec application côté serveur (JWT claims) et restrictions de navigation côté client.
  \item \textbf{Centraliser les logs} : mettre en place un pipeline automatisé (Filebeat -> Logstash -> Elasticsearch) collectant et enrichissant tous les journaux applicatifs dans un moteur de recherche unifié, avec rétention automatique ILM à 14 jours.
  \item \textbf{Détecter les incidents} : implémenter une détection automatique des événements de sécurité : brute force, accès interdits, rate limiting, uploads suspects, injections SQL, User-Agents offensifs.
  \item \textbf{Alertes en temps réel (SSE)} : diffuser les événements de sécurité critiques vers le dashboard en temps réel via Server-Sent Events, sans nécessiter de polling.
  \item \textbf{Monitorer en temps réel} : déployer un système de métriques (Prometheus), de visualisation (Grafana) et d'alertes (Alertmanager avec webhook) couvrant l'infrastructure et la sécurité.
  \item \textbf{Offrir une interface unifiée} : proposer un tableau de bord interactif centralisant logs, alertes, métriques, géolocalisation et état des services, avec sélecteur de plage temporelle, export CSV et gestion des thèmes.
  \item \textbf{Démontrer la résilience} : simuler des scénarios d'attaque réels pour valider l'efficacité de l'ensemble du dispositif de sécurité.
\end{enumerate}
\end{successbox}

\section{Planification}

Le projet a été réalisé sur une durée de six semaines, avec un chevauchement intentionnel des phases permettant une intégration progressive des composants. Le tableau~\ref{tab:gantt} présente la répartition des tâches.

\begin{table}[H]
  \centering
  \small
  \setlength{\tabcolsep}{4pt}
  \caption{Diagramme de Gantt - Planification du projet (6 semaines)}
  \label{tab:gantt}
  \begin{tabular}{|l|c|c|c|c|c|c|}
    \hline
    \rowcolor{NavyBlue!15}
    \textbf{Phase / Tâche} & \textbf{S1} & \textbf{S2} & \textbf{S3} & \textbf{S4} & \textbf{S5} & \textbf{S6} \\
    \hline
    Analyse des besoins et état de l'art
      & \cellcolor{emsigreen!60} & \cellcolor{emsigreen!60} & & & & \\
    \hline
    Conception de l'architecture Docker
      & \cellcolor{emsigreen!60} & \cellcolor{emsigreen!60} & & & & \\
    \hline
    Déploiement Docker Compose et réseaux
      & & \cellcolor{emsigreen!60} & \cellcolor{emsigreen!60} & & & \\
    \hline
    Configuration API Gateway (TLS, rate limiting)
      & & \cellcolor{emsigreen!60} & \cellcolor{emsigreen!60} & & & \\
    \hline
    Développement microservices Flask
      & & & \cellcolor{emsigreen!60} & \cellcolor{emsigreen!60} & & \\
    \hline
    Mise en place pipeline ELK Stack
      & & & \cellcolor{emsigreen!60} & \cellcolor{emsigreen!60} & & \\
    \hline
    Monitoring Prometheus / Grafana / Alertmanager
      & & & & \cellcolor{emsigreen!60} & \cellcolor{emsigreen!60} & \\
    \hline
    Dashboard SecureWatch (frontend)
      & & & & & \cellcolor{emsigreen!60} & \cellcolor{emsigreen!60} \\
    \hline
    Tests, simulation d'attaques et validation
      & & & & & \cellcolor{emsigreen!60} & \cellcolor{emsigreen!60} \\
    \hline
    Rédaction du rapport
      & & & & & \cellcolor{emsigreen!30} & \cellcolor{emsigreen!60} \\
    \hline
  \end{tabular}
\end{table}

\medskip
Nous avons découpé le projet en six semaines en adoptant une approche incrémentale : chaque couche de l'architecture est construite et testée avant d'en intégrer la suivante. Les deux premières semaines ont été consacrées à la conception — comprendre les besoins, choisir les outils, dessiner l'architecture. À partir de la semaine 2, l'implémentation démarre : gateway, microservices, pipeline ELK, monitoring, puis le dashboard en dernière étape. Le chevauchement intentionnel de certaines phases permet d'identifier rapidement les problèmes d'intégration. Les tests et la rédaction du rapport occupent les deux dernières semaines en parallèle.

% ============================================================
%  CHAPITRE 2 * ÉTAT DE L'ART
% ============================================================
\chapter{État de l'Art et Fondements Théoriques}

\section{Introduction}

Avant de construire quoi que ce soit, il faut comprendre ce qui existe déjà et pourquoi certaines technologies s'imposent naturellement. Ce chapitre explore les concepts fondamentaux derrière SecureWatch — microservices, sécurité distribuée, centralisation de logs — et justifie, comparaison à l'appui, chaque outil retenu.

\section{Architecture Microservices et Conteneurisation}

\subsection{Le paradigme microservices}

L'idée derrière les microservices est simple : plutôt qu'un gros bloc monolithique, on découpe l'application en petits services indépendants, chacun responsable d'une fonctionnalité précise. Ces services communiquent entre eux via des APIs bien définies, généralement en HTTP/REST ou via des files de messages.

Les avantages sont réels :
\begin{itemize}[leftmargin=*]
  \item \textbf{Déploiement indépendant} : mettre à jour un service ne touche pas les autres.
  \item \textbf{Scalabilité granulaire} : on réplique uniquement les services sous pression, pas toute l'application.
  \item \textbf{Isolation des pannes} : la défaillance d'un service ne fait pas tomber l'ensemble du système.
  \item \textbf{Liberté technologique} : chaque service peut utiliser le langage ou la stack les mieux adaptés à son rôle.
\end{itemize}

Mais cette approche a un revers : la complexité opérationnelle explose. La communication réseau entre services, la cohérence des données, la sécurisation de chaque point d'accès, et la nécessité d'avoir une vision globale du système — tout cela demande des outils spécifiques qu'un environnement monolithique ne nécessite pas.

\subsection{Docker et la conteneurisation}

\textbf{Docker} a transformé la conteneurisation en pratique courante. Un conteneur encapsule une application avec toutes ses dépendances dans un environnement isolé et reproductible. Plus légers que les machines virtuelles (ils partagent le noyau de l'OS hôte), ils démarrent en quelques secondes et s'exécutent de manière identique sur n'importe quelle machine.

\textbf{Docker Compose} pousse l'idée plus loin pour les environnements locaux. Un simple fichier \texttt{docker-compose.yml} suffit à décrire toute une infrastructure — services, réseaux, volumes, variables d'environnement — et une seule commande (\texttt{docker-compose up}) lance le tout. Pour un projet comme SecureWatch avec ses 13 services, c'est indispensable.

\section{La Cybersécurité en Environnement Distribué}

\subsection{API Gateway et routage sécurisé}

Dans une architecture microservices, exposer directement chaque service sur internet serait une catastrophe sécuritaire. C'est là qu'intervient l'\textbf{API Gateway} : un point d'entrée unique qui filtre tout le trafic entrant avant de le router vers les services appropriés. En centralisant les fonctions de sécurité à cet endroit, on évite de les implémenter dans chaque service :

\begin{itemize}[leftmargin=*]
  \item \textbf{Terminaison TLS} : déchiffrement des connexions HTTPS au niveau de la gateway, évitant que chaque microservice ne gère son propre certificat.
  \item \textbf{Rate Limiting} : limitation du débit des requêtes par unité de temps et par IP, protection contre les attaques par force brute et les abus d'API.
  \item \textbf{Authentification} : vérification des tokens ou des credentials avant transmission de la requête au service destinataire.
  \item \textbf{Routage intelligent} : redirection des requêtes vers les microservices appropriés selon des règles basées sur le chemin, les en-têtes ou d'autres critères.
  \item \textbf{Journalisation centralisée} : enregistrement structuré de toutes les requêtes entrantes, constituant la première couche de détection d'incidents.
\end{itemize}

Le rate limiting repose souvent sur l'algorithme du \textit{leaky bucket} (seau percé). L'image est parlante : les requêtes arrivent à un rythme variable, comme de l'eau versée dans un seau, mais s'écoulent à débit constant. Dès que le seau déborde, les nouvelles requêtes sont rejetées avec un code HTTP 429. Simple, efficace, et particulièrement redoutable contre les attaques par dictionnaire qui misent sur la vitesse.

\subsection{Gestion des identités et protocoles d'authentification}

\textbf{OAuth 2.0} est le standard d'autorisation qui régit aujourd'hui la plupart des échanges sécurisés entre services. Il définit plusieurs flux selon le contexte : Authorization Code pour les applications web, Client Credentials pour les communications service-à-service. Il ne s'occupe pas de l'authentification elle-même, mais de déléguer l'accès à des ressources protégées.

Les \textbf{JSON Web Tokens (JWT)} sont devenus le format d'échange de tokens incontournable en microservices. Un JWT est un objet en trois parties (header, payload, signature) encodées en Base64 et séparées par des points. Son atout principal est son caractère \textit{stateless} : pas besoin de consulter une base de données pour valider un token — la signature suffit. Ce que le token contient (identité, rôles, expiration) est vérifiable sans état côté serveur.

Le \textbf{RBAC} (Role-Based Access Control) complète ce tableau : plutôt que d'attribuer des permissions à chaque utilisateur individuellement, on les associe à des rôles (admin, user, operator). Un utilisateur hérite des permissions de son rôle. Simple, lisible, facile à maintenir.

\subsection{Surveillance et détection d'incidents (SIEM)}

Un \textbf{SIEM} (Security Information and Event Management) est en quelque sorte la tour de contrôle d'une infrastructure sécurisée. Il agrège les événements venant de toutes les sources — serveurs web, pare-feu, applications, bases de données — et les analyse pour détecter des comportements anormaux. Ses fonctions fondamentales :

\begin{itemize}[leftmargin=*]
  \item \textbf{Collecte et normalisation} : rassembler des logs hétérogènes dans un format commun et exploitable.
  \item \textbf{Corrélation d'événements} : identifier des patterns suspects qui ne sont visibles qu'en croisant plusieurs sources.
  \item \textbf{Alertes} : notifier automatiquement les équipes de sécurité dès qu'un comportement suspect est détecté.
  \item \textbf{Archivage} : conserver les preuves pour l'investigation forensique et les obligations de conformité.
\end{itemize}

Pour SecureWatch, nous construisons un SIEM léger autour de l'ELK Stack. Pas aussi complet qu'un Splunk d'entreprise, mais fonctionnel, gratuit, et déployable en local — exactement ce qu'il nous faut dans ce contexte.

\section{Technologies Retenues et Justification}

\begin{table}[H]
  \centering
  \caption{Technologies utilisées dans SecureWatch et justification des choix}
  \begin{tabularx}{\textwidth}{@{}llX@{}}
    \toprule
    \textbf{Technologie} & \textbf{Version} & \textbf{Justification du choix} \\
    \midrule
    Nginx               & Alpine  & Légèreté, performance élevée, modules SSL et rate limiting natifs, images Docker officielles Alpine très compactes (\textasciitilde25 Mo). \\
    \addlinespace
    Elasticsearch       & 8.12.0  & Moteur de recherche distribué, full-text, scalable. Standard de facto pour la centralisation de logs. Supporte les requêtes d'agrégation complexes nécessaires au dashboard. \\
    \addlinespace
    Logstash            & 8.12.0  & Pipeline ETL flexible avec filtres Grok et Ruby. Permet l'enrichissement des logs avec des tags de sécurité personnalisés avant indexation. \\
    \addlinespace
    Kibana              & 8.12.0  & Visualisation native ELK, tableaux de bord avancés, interface Discover pour l'investigation. La cohérence de version 8.12.0 avec ES et Logstash est obligatoire. \\
    \addlinespace
    Filebeat            & 8.12.0  & Agent léger de collecte de logs (faible empreinte CPU/RAM). Surveille les fichiers en temps réel et retransmet les nouvelles lignes à Logstash. \\
    \addlinespace
    Prometheus          & Latest  & Standard de facto pour les métriques cloud-native. Modèle pull, TSDB intégré, langage PromQL puissant pour les alertes. \\
    \addlinespace
    Grafana             & Latest  & Dashboards riches, auto-provisioning via fichiers YAML/JSON. Compatible avec Prometheus nativement. \\
    \addlinespace
    Alertmanager        & Latest  & Routage d'alertes, déduplication, inhibition et groupement. S'intègre nativement avec Prometheus. \\
    \addlinespace
    Python / Flask      & 3.11    & Rapidité de développement, bibliothèque \texttt{prometheus-flask-exporter} native pour l'instrumentation. Légèreté des images Docker. \\
    \addlinespace
    Docker Compose      & Latest  & Orchestration locale simple, fichier YAML déclaratif, support des réseaux et volumes. Reproductibilité totale de l'environnement. \\
    \bottomrule
  \end{tabularx}
\end{table}

\medskip
Ces choix ne sont pas arbitraires. Chaque outil retenu satisfait plusieurs critères : images Docker officielles disponibles, communauté active, documentation solide, et compatibilité avec les autres composants. Un point important à noter : la suite ELK impose une cohérence stricte de version — Elasticsearch, Logstash, Kibana et Filebeat doivent être sur la même version (8.12.0 ici), sans quoi des incompatibilités de protocole apparaissent. Prometheus et Grafana sont plus souples et utilisent la dernière version stable.

Deux choix méritent une justification particulière. \textbf{Nginx plutôt que Kong} : Kong est fonctionnellement plus riche, mais il traîne avec lui une base de données PostgreSQL et une configuration nettement plus complexe. Pour un projet en environnement local, ce surcoût opérationnel n'est pas justifié. \textbf{Flask personnalisé plutôt que Keycloak} : un Identity Provider complet aurait masqué la logique de détection que nous voulions implémenter et observer. Flask nous donne un contrôle total sur le comportement d'authentification, ce qui est essentiel pour la simulation.

\section{Étude Comparative}

\subsection{Comparaison des solutions d'API Gateway}

\begin{table}[H]
  \centering
  \caption{Comparaison des principales solutions d'API Gateway}
  \begin{tabularx}{\textwidth}{@{}lXXXXX@{}}
    \toprule
    \textbf{Solution} & \textbf{Licence} & \textbf{Rate Limiting} & \textbf{Auth intégrée} & \textbf{Config} & \textbf{Retenu} \\
    \midrule
    \textbf{Nginx}   & Open Source & Natif (zones) & Basique & Fichier texte & \cmark \\
    Kong             & Apache 2.0  & Plugin        & OAuth2, JWT, OIDC & Admin API / YAML & \xmark \\
    Traefik          & MIT         & Middleware    & Basique           & Labels Docker & \xmark \\
    HAProxy          & LGPL        & Limité        & Non               & Fichier texte & \xmark \\
    \bottomrule
  \end{tabularx}
\end{table}

\medskip
Le tableau parle de lui-même. Nginx s'impose pour sa simplicité : rate limiting natif, zéro dépendance externe, configuration lisible en quelques dizaines de lignes. Kong est plus puissant, mais il traîne avec lui une base de données PostgreSQL qu'on ne veut pas gérer en local. Traefik est excellent en Kubernetes, moins à l'aise en Docker Compose pur. HAProxy fait du load balancing, pas vraiment de la sécurité applicative.

\subsection{Comparaison des solutions de gestion de logs}

\begin{table}[H]
  \centering
  \caption{Comparaison des solutions de centralisation de logs}
  \begin{tabularx}{\textwidth}{@{}lXXXXX@{}}
    \toprule
    \textbf{Solution} & \textbf{Type} & \textbf{Licence} & \textbf{Alertes} & \textbf{Scalabilité} & \textbf{Retenu} \\
    \midrule
    \textbf{ELK Stack} & Distribué & Open Source & Via Watcher/ES & Très élevée & \cmark \\
    Splunk             & Centralisé & Commercial  & Native, avancée  & Très élevée & \xmark \\
    Graylog            & Distribué  & Open Source & Native           & Élevée      & \xmark \\
    Loki + Grafana     & Centralisé & Open Source & Via Grafana      & Élevée      & \xmark \\
    \bottomrule
  \end{tabularx}
\end{table}

\medskip
ELK s'impose comme le choix naturel : l'écosystème est complet, la documentation abondante, les images Docker officielles disponibles. Splunk ferait mieux techniquement, mais son modèle de facturation (par volume ingéré) est incompatible avec un projet pédagogique. Graylog est une alternative sérieuse, mais sa communauté plus réduite représente un risque si on bloque sur un problème. Loki est très léger et s'intègre bien avec Grafana, mais son architecture sans indexation du contenu des logs (seulement des labels) limiterait nos capacités de recherche full-text dans le Dashboard.

\subsection{Comparaison des solutions de monitoring}

\begin{table}[H]
  \centering
  \caption{Comparaison des solutions de monitoring}
  \begin{tabularx}{\textwidth}{@{}lXXXXX@{}}
    \toprule
    \textbf{Solution} & \textbf{Modèle} & \textbf{Licence} & \textbf{Alertes} & \textbf{Dashboards} & \textbf{Retenu} \\
    \midrule
    \textbf{Prometheus + Grafana} & Pull & Open Source & Alertmanager & Grafana & \cmark \\
    InfluxDB + Telegraf           & Push & Open Source & Kapacitor    & Chronograf & \xmark \\
    Datadog                       & Agent & SaaS        & Native       & Native & \xmark \\
    Zabbix                        & Agent & Open Source & Native       & Native & \xmark \\
    \bottomrule
  \end{tabularx}
\end{table}

\medskip
Prometheus + Grafana + Alertmanager, c'est le trio standard du monitoring cloud-native, et pour de bonnes raisons. Le modèle pull de Prometheus (il va chercher les métriques plutôt que les recevoir) est parfait pour Docker : les endpoints \texttt{/metrics} sont accessibles en réseau interne, pas besoin d'ouvrir des ports supplémentaires. Grafana transforme ces données en dashboards riches, et Alertmanager gère le routage des alertes avec déduplication. Datadog est plus avancé, mais sa facturation par hôte (SaaS) n'est pas envisageable dans un contexte pédagogique local.

% ============================================================
%  CHAPITRE 3 * MISE EN ŒUVRE TECHNIQUE
% ============================================================
\chapter{Mise en Œuvre Technique}

\section{Introduction}

Passons maintenant à l'implémentation concrète. Tout tourne sur un poste Windows 11 avec Docker Desktop — pas de serveur distant, pas de cloud, juste une machine de développement. Ce chapitre décrit comment les 13 services ont été assemblés, configurés et sécurisés, de la gateway en entrée jusqu'aux dashboards de surveillance.

\section{Déploiement de l'Infrastructure}

\subsection{Architecture globale}

SecureWatch repose sur une architecture microservices containerisée via Docker Compose, composée de \textbf{14 services} (dont un conteneur d'initialisation \texttt{es-setup} à démarrage unique) répartis sur \textbf{3 réseaux Docker isolés}. L'architecture suit le principe de \textit{défense en profondeur} : chaque couche ajoute une barrière de sécurité supplémentaire.

\screenshot{architecture_diagram}{Architecture globale de SecureWatch - 13 services sur 3 réseaux Docker isolés}{architecture}

\medskip
L'architecture s'organise en trois couches distinctes. En façade, la gateway Nginx : c'est le seul point de contact avec l'extérieur. En zone interne, les microservices, totalement invisibles depuis l'hôte. En arrière-plan, les systèmes de surveillance qui observent tout sans interférer avec le trafic applicatif. Deux flux traversent cette architecture : les requêtes des utilisateurs, et la télémétrie (logs et métriques) que les services génèrent en permanence.

\begin{infobox}[Principe fondamental d'isolation réseau]
Les microservices (auth-service et api-service) n'ont \textbf{aucun port exposé} sur l'hôte. Tout trafic externe passe obligatoirement par l'API Gateway sur les ports 8080 (HTTP -> redirection) et 8443 (HTTPS). Un attaquant ne peut pas atteindre directement les services internes même s'il connaît leurs adresses IP Docker.
\end{infobox}

\subsection{Configuration des réseaux Docker}

L'isolation réseau est assurée par trois réseaux Docker bridge distincts, chacun dédié à un flux de données spécifique :

\begin{table}[H]
  \centering
  \caption{Réseaux Docker et leurs rôles dans l'architecture SecureWatch}
  \begin{tabularx}{\textwidth}{@{}lX@{}}
    \toprule
    \textbf{Réseau} & \textbf{Services inclus \& rôle} \\
    \midrule
    \texttt{gateway-net}
      & \textit{gateway, auth-service, api-service, redis, nginx-exporter, frontend.}
        Réseau de routage applicatif : la gateway est le seul point d'entrée vers les microservices. Redis y est rattaché pour que auth-service et api-service y accèdent directement pour la persistance des compteurs. \\
    \addlinespace
    \texttt{monitoring-net}
      & \textit{prometheus, alertmanager, grafana, node-exporter, nginx-exporter, auth-service, api-service, gateway.}
        Canal de collecte de métriques interne. Prometheus scrape les endpoints \texttt{/metrics} de tous les services toutes les 15 secondes. \\
    \addlinespace
    \texttt{elk}
      & \textit{elasticsearch, logstash, kibana, filebeat, auth-service, api-service, gateway, frontend.}
        Canal de centralisation des logs. Filebeat envoie les journaux vers Logstash, qui les indexe dans Elasticsearch. \\
    \bottomrule
  \end{tabularx}
\end{table}

\medskip
Ce cloisonnement n'est pas qu'organisationnel — il est sécuritaire. Dans Docker, deux conteneurs sur des réseaux différents ne peuvent pas communiquer, même s'ils tournent sur la même machine. Concrètement : si un attaquant compromet Kibana, il se retrouve coincé sur le réseau \texttt{elk}, sans accès aux microservices sur \texttt{gateway-net}. La compromission d'un composant de surveillance ne donne pas accès à l'application — c'est l'essence de la défense en profondeur.

\textbf{Redis — persistance des compteurs de brute force.} Un conteneur Redis 7 Alpine est ajouté sur \texttt{gateway-net} avec une limite mémoire de 32 Mo et la politique d'éviction \texttt{allkeys-lru}. Les compteurs de tentatives d'authentification échouées par IP y sont stockés avec une fenêtre glissante d'une heure (\texttt{EXPIRE 3600}). Sans Redis, ces compteurs disparaissaient à chaque redémarrage du conteneur auth-service, permettant à un attaquant de réinitialiser son compteur en forçant un redémarrage. Avec Redis, les compteurs survivent aux redémarrages. auth-service dispose d'un fallback automatique vers la mémoire locale si Redis est injoignable — le service ne tombe jamais à cause de Redis.

\textbf{ILM — rétention automatique des index Elasticsearch.} Un conteneur d'initialisation \texttt{es-setup} (one-shot, s'exécute une fois puis s'arrête) applique au démarrage la politique ILM \texttt{securewatch-14d} sur Elasticsearch. Cette politique supprime automatiquement les index \texttt{security-logs-*} de plus de 14 jours. Sans ILM, les index s'accumulent indéfiniment sur le disque. Un template d'index associe automatiquement tout nouvel index \texttt{security-logs-*} à cette politique, sans intervention manuelle.

\subsection{Flux de données}

\screenshot[0.90\textwidth]{flux_donnees}{Flux de données entre les composants - logs (ELK) et métriques (Prometheus)}{flux}

\medskip
Deux flux parallèles alimentent le système de surveillance. Le \textbf{flux de logs} suit la chaîne ELK classique : chaque événement est écrit en JSON dans un fichier sur un volume partagé, Filebeat le récupère en temps quasi-réel, Logstash l'enrichit avec des tags de sécurité avant de l'indexer dans Elasticsearch. Kibana et le Dashboard peuvent ensuite interroger ces données librement. Le \textbf{flux de métriques} fonctionne différemment : chaque service publie un endpoint \texttt{/metrics} que Prometheus vient lire toutes les 15 secondes. Il évalue ses règles d'alertes toutes les 30 secondes et, quand une règle se déclenche, notifie Alertmanager. Grafana visualise l'ensemble en temps réel.

\subsection{Mise en place de l'API Gateway (Nginx)}

L'API Gateway Nginx constitue l'unique point d'entrée vers les microservices. Elle écoute sur deux ports et remplit quatre fonctions essentielles :

\begin{enumerate}[leftmargin=*]
  \item \textbf{Terminaison TLS} : déchiffrement HTTPS, certificat auto-signé RSA 2048 bits.
  \item \textbf{Rate limiting} : protection contre le brute force et les abus d'API via l'algorithme du leaky bucket.
  \item \textbf{Reverse proxy} : routage des requêtes vers les microservices internes par préfixe de chemin.
  \item \textbf{Journalisation JSON} : logs structurés de chaque requête avec IP source, méthode, URI, code de réponse et temps de traitement.
\end{enumerate}

\section{Configuration de la Sécurité}

\subsection{Terminaison TLS et en-têtes de sécurité}

La gateway écoute sur deux ports avec des comportements distincts :

\begin{table}[H]
  \centering
  \caption{Ports de l'API Gateway et comportements associés}
  \begin{tabular}{@{}lll@{}}
    \toprule
    \textbf{Port} & \textbf{Protocole} & \textbf{Comportement} \\
    \midrule
    8080 & HTTP  & Redirection 301 permanente vers \texttt{https://<host>:8443} \\
    8443 & HTTPS & Toutes les routes API (\texttt{/auth/*}, \texttt{/api/*}, \texttt{/health}) \\
    \bottomrule
  \end{tabular}
\end{table}

\medskip
Le port 8080 existe uniquement pour une raison : rediriger immédiatement vers HTTPS. Personne ne doit pouvoir communiquer en clair avec les services. Le certificat utilisé est auto-signé (RSA 2048 bits, généré au démarrage du conteneur) — acceptable en développement, à remplacer par Let's Encrypt ou un certificat d'entreprise en production.

En plus du chiffrement, la gateway applique systématiquement un ensemble d'en-têtes de sécurité recommandés par l'OWASP :

\begin{lstlisting}[style=bash, caption={En-têtes de sécurité HTTP appliqués par la gateway}]
Strict-Transport-Security: max-age=31536000; includeSubDomains
X-Frame-Options: DENY
X-Content-Type-Options: nosniff
X-XSS-Protection: 1; mode=block
X-Gateway: SecureWatch-PFA-2026
\end{lstlisting}

\texttt{HSTS} oblige le navigateur à utiliser HTTPS pour toute requête future vers ce domaine, pendant un an. \texttt{X-Frame-Options: DENY} bloque l'intégration en iframe (protection contre le clickjacking). \texttt{X-Content-Type-Options: nosniff} empêche le navigateur de deviner le type MIME d'une réponse — une précaution simple mais efficace contre certaines variantes d'attaques XSS.

\subsection{Rate Limiting - Protection contre le brute force et les abus}

Nginx implémente le rate limiting via l'algorithme du \textit{leaky bucket}. Trois zones sont définies, chacune ciblant un type de trafic différent :

\begin{table}[H]
  \centering
  \caption{Zones de rate limiting configurées dans l'API Gateway}
  \begin{tabular}{@{}llll@{}}
    \toprule
    \textbf{Zone} & \textbf{Limite} & \textbf{Burst} & \textbf{Cible} \\
    \midrule
    \texttt{auth\_zone}   & 5 req/min par IP  & 10  & Toutes les routes \texttt{/auth/*} \\
    \texttt{api\_zone}    & 30 req/min par IP & 50  & Toutes les routes \texttt{/api/*} \\
    \texttt{global\_zone} & 100 req/min par IP & 200 & Protection DDoS globale \\
    \bottomrule
  \end{tabular}
\end{table}

\medskip
Les trois zones ont des seuils calibrés pour leurs usages respectifs. L'\texttt{auth\_zone} est la plus stricte : 5 requêtes par minute, c'est intentionnellement bas. Un humain qui saisit son mot de passe n'a pas besoin de plus. Un script d'attaque qui essaie 100 mots de passe à la seconde, lui, sera bloqué presque immédiatement. Le paramètre \texttt{burst=10} sert à absorber des rafales légitimes courtes (rechargement de page, appels parallèles) sans déclencher de faux positifs. L'\texttt{api\_zone} est plus permissive pour un usage normal. La \texttt{global\_zone} est le dernier filet — protection contre un flood non ciblé.

Quand la limite est dépassée, la gateway répond avec HTTP 429 et un message JSON, tout en loggant l'événement dans ELK avec le tag \texttt{rate\_limited} :

\begin{lstlisting}[style=yaml, caption={Réponse HTTP 429 personnalisée retournée par la gateway}]
{
  "error": "Too Many Requests",
  "message": "Rate limit exceeded. Please slow down.",
  "service": "api-gateway"
}
\end{lstlisting}

\medskip
La structure de cette réponse est volontaire : elle indique clairement la nature du blocage (\texttt{error}), fournit un message lisible par un humain (\texttt{message}), et identifie la source du blocage (\texttt{service: api-gateway}). En parallèle, la gateway enregistre l'IP source, la route tentée et le code 429 dans ses logs JSON — ces entrées remontent dans Elasticsearch via le pipeline ELK et alimentent les compteurs du Dashboard.

\screenshot[0.85\textwidth]{gateway_rate_limit}{Démonstration du rate limiting - terminal lors du scénario rate-limit}{ratelimit}

\medskip
Le terminal ci-dessus affiche la sortie du script de simulation lors du scénario de test rate-limit. Les premières requêtes reçoivent le code HTTP 200 (succès) tandis que celles dépassant le burst configuré retournent HTTP 429. La progression des codes de retour confirme le bon fonctionnement de l'algorithme leaky bucket : les requêtes initiales passent, puis le rejet devient systématique une fois le seuil dépassé.

\subsection{Microservices Flask - Authentification et Détection}

\subsubsection{Auth Service (port interne 5001)}

Le service d'authentification gère les connexions et émet des logs de sécurité pour chaque tentative. Il expose également des métriques Prometheus via \texttt{prometheus-flask-exporter}.

\begin{table}[H]
  \centering
  \caption{Endpoints exposés par auth-service}
  \begin{tabular}{@{}lll@{}}
    \toprule
    \textbf{Méthode} & \textbf{Route} & \textbf{Description} \\
    \midrule
    GET  & \texttt{/health}   & Health check - accessible via gateway à \texttt{/auth/health} \\
    POST & \texttt{/login}    & Authentification - retourne un token si succès (HTTP 200) ou échec (HTTP 401) \\
    POST & \texttt{/logout}   & Déconnexion \\
    POST & \texttt{/register} & Enregistrement d'un nouvel utilisateur \\
    GET  & \texttt{/metrics}  & Métriques Prometheus scrappées par Prometheus toutes les 15 s \\
    \bottomrule
  \end{tabular}
\end{table}

\medskip
Chaque endpoint d'authentification génère des métriques Prometheus spécifiques pour le suivi des indicateurs de sécurité :

\begin{lstlisting}[language=Python, caption={Compteurs Prometheus définis dans auth-service pour la détection de sécurité}]
from prometheus_flask_exporter import PrometheusMetrics
from prometheus_client import Counter

metrics = PrometheusMetrics(app, default_labels={'service': 'auth-service'})

auth_failures_total = Counter(
    'auth_failures_total',
    'Total authentication failures',
    ['ip']
)
brute_force_total = Counter(
    'brute_force_total',
    'Brute force attacks detected',
    ['ip']
)
\end{lstlisting}

La \textbf{détection brute force} est implémentée directement dans auth-service : si une même IP effectue \textbf{5 échecs d'authentification ou plus} en peu de temps, un événement \texttt{brute\_force} est loggé au niveau CRITICAL et le compteur \texttt{brute\_force\_total} est incrémenté — déclenchant ainsi les alertes Prometheus. Les compteurs sont désormais persistés dans Redis (fenêtre glissante de 1 heure) et survivent aux redémarrages de conteneur.

Auth-service intègre également deux mécanismes de détection supplémentaires actifs sur toutes les requêtes via un middleware \texttt{@before\_request} :

\begin{itemize}[leftmargin=*]
  \item \textbf{Détection d'injection SQL} : le champ \texttt{username} est analysé par une expression régulière couvrant les patterns classiques (\texttt{UNION}, \texttt{SELECT}, \texttt{OR 1=1}, opérateurs hexadécimaux...). Toute correspondance déclenche un log \texttt{sqli\_attempt} et retourne HTTP 400.
  \item \textbf{Détection de User-Agents offensifs} : les requêtes dont le \texttt{User-Agent} correspond à des outils d'attaque connus (sqlmap, nikto, nmap, gobuster, ffuf, nuclei, burpsuite...) génèrent un log \texttt{suspicious\_ua} et incrémentent le compteur \texttt{scanner\_ua\_total}.
\end{itemize}

\textbf{SSE — alertes en temps réel.} Auth-service expose un endpoint \texttt{/events/stream} implémentant les Server-Sent Events. Le frontend s'y connecte via \texttt{EventSource} dès l'authentification (le JWT est passé en paramètre \texttt{?token=} puisque l'API EventSource ne permet pas d'envoyer de headers personnalisés). Dès qu'un événement \texttt{auth\_failure} ou \texttt{brute\_force} se produit, il est poussé en temps réel à tous les clients connectés. Côté Dashboard, un événement brute force déclenche immédiatement une notification toast violette et un rafraîchissement des compteurs, sans attendre le prochain cycle de polling de 8 secondes.

\screenshot[0.90\textwidth]{ui_review/auth_login_screen}{Écran de connexion SecureWatch — authentification JWT avant accès au dashboard}{login_screen}

Les comptes de test disponibles pour la simulation sont détaillés dans le tableau suivant :

\begin{table}[H]
  \centering
  \caption{Comptes utilisateurs de test disponibles dans auth-service}
  \begin{tabular}{@{}lll@{}}
    \toprule
    \textbf{Utilisateur} & \textbf{Mot de passe} & \textbf{Rôle} \\
    \midrule
    \texttt{admin}    & \texttt{password123} & Administrateur \\
    \texttt{user1}    & \texttt{secret456}   & Utilisateur standard \\
    \texttt{operator} & \texttt{ops789}      & Opérateur \\
    \texttt{hicham}   & \texttt{pfa2026}     & Utilisateur standard \\
    \bottomrule
  \end{tabular}
\end{table}

\medskip
Ces comptes sont codés en dur — c'est volontaire et acceptable dans le cadre d'une simulation pédagogique. En production, il n'en serait évidemment pas question : les mots de passe seraient hashés (bcrypt), stockés en base, et l'authentification déléguée à un Identity Provider dédié (Keycloak, Auth0).

\subsubsection{API Service (port interne 5002)}

Le service API expose les ressources de l'application avec une application côté serveur du RBAC. Chaque endpoint protégé vérifie la signature JWT \textit{et} le rôle de l'utilisateur avant de répondre — une architecture à deux niveaux de vérification :

\begin{table}[H]
  \centering
  \caption{Endpoints de api-service et comportements de sécurité associés}
  \begin{tabular}{@{}lll@{}}
    \toprule
    \textbf{Méthode} & \textbf{Route} & \textbf{Résultat / Événement généré} \\
    \midrule
    GET  & \texttt{/api/data}    & Tous les rôles authentifiés — HTTP 200 (\texttt{api\_access}) ou 401 \\
    GET  & \texttt{/api/users}   & Admin + Operator uniquement — 200 ou 401/403 selon rôle \\
    GET  & \texttt{/api/reports} & Admin + Operator uniquement — 200 ou 401/403 selon rôle \\
    GET  & \texttt{/api/admin}   & Zone bloquée — HTTP 403 (\texttt{forbidden\_access}) pour tous \\
    GET  & \texttt{/api/config}  & Zone bloquée — HTTP 403 (\texttt{forbidden\_access}) pour tous \\
    POST & \texttt{/api/upload}  & Admin + Operator — 200 ou 415 (\texttt{suspicious\_upload}), ou 401/403 \\
    POST & \texttt{/api/alerts/webhook} & Récepteur Alertmanager (réseau interne) — stockage Redis \\
    GET  & \texttt{/api/alerts/pushed}  & Admin + Operator — alertes reçues via webhook (Redis) \\
    GET  & \texttt{/health}      & Health check \\
    GET  & \texttt{/metrics}     & Métriques Prometheus \\
    \bottomrule
  \end{tabular}
\end{table}

La fonction \texttt{require\_roles()} extrait le claim \texttt{role} du JWT signé et vérifie l'appartenance à la liste des rôles autorisés. Un token absent retourne HTTP 401 ; un token valide mais au rôle insuffisant retourne HTTP 403 et génère un log \texttt{forbidden\_access}. Ce comportement différencié permet au Dashboard de distinguer un accès non authentifié d'un accès refusé pour raison de rôle.

Api-service intègre également un middleware \texttt{detect\_threats()} actif sur toutes les requêtes :

\begin{table}[H]
  \centering
  \caption{Mécanismes de détection actifs dans api-service (middleware \texttt{@before\_request})}
  \begin{tabularx}{\textwidth}{@{}lXl@{}}
    \toprule
    \textbf{Détection} & \textbf{Mécanisme} & \textbf{Compteur Prometheus} \\
    \midrule
    User-Agent offensif & Regex sur sqlmap, nikto, nmap, gobuster, ffuf, nuclei, burpsuite... & \texttt{scanner\_ua\_total} \\
    Anomalie géographique & Correspondance de préfixe IP (Tor, Russie, Chine, Corée du Nord, Iran, Nigeria) & \texttt{geo\_anomaly\_total} \\
    \bottomrule
  \end{tabularx}
\end{table}

\medskip
La détection géographique analyse le premier hop du header \texttt{X-Forwarded-For} (l'IP déclarée par le client) contre une liste de préfixes connus comme sources d'activité malveillante. Cela inclut les nœuds de sortie Tor (\texttt{185.220.*}, \texttt{162.247.*}), plusieurs plages allouées à des pays à fort volume d'attaques et des opérateurs cloud fréquemment utilisés pour des campagnes automatisées. Une détection génère un log \texttt{geo\_anomaly} avec le code pays et le nom de la région, visible dans Kibana et le Dashboard.

\textbf{Webhook Alertmanager.} L'endpoint \texttt{/api/alerts/webhook} reçoit les notifications POST d'Alertmanager (configuré comme receiver \texttt{webhook-api}). Chaque alerte reçue est sérialisée et stockée dans Redis (\texttt{LPUSH webhook\_alerts}, capped à 100 entrées). L'endpoint \texttt{/api/alerts/pushed} expose ces alertes aux rôles admin et operator, permettant au Dashboard d'afficher les alertes Prometheus directement sans passer par l'interface Alertmanager.

\section{Surveillance et Audit}

\subsection{Pipeline de Logs ELK Stack}

\subsubsection{Vue d'ensemble}

Le pipeline de centralisation des logs suit la chaîne :

\begin{center}
  \texttt{Services} $\rightarrow$ \texttt{/logs/*.log} $\rightarrow$
  \texttt{Filebeat} $\rightarrow$ \texttt{Logstash:5044} $\rightarrow$
  \texttt{Elasticsearch:9200} $\rightarrow$ \texttt{Kibana / Dashboard}
\end{center}

Chaque service écrit ses logs au format JSON structuré dans un volume Docker partagé (\texttt{/logs/}). Ce format structuré est essentiel pour que Logstash puisse parser et enrichir les événements sans recourir à des règles Grok complexes.

\begin{lstlisting}[style=yaml, caption={Exemple de log brute\_force généré par auth-service * format JSON structuré}]
{
  "timestamp": "2026-05-16T02:54:36.894Z",
  "service": "auth-service",
  "level": "CRITICAL",
  "event_type": "brute_force",
  "ip": "185.220.101.5",
  "failed_attempts": 10,
  "alert_type": "BRUTE_FORCE",
  "message": "BRUTE FORCE ATTACK detected from IP 185.220.101.5"
}
\end{lstlisting}

\medskip
Cet exemple est tiré d'une vraie session de simulation. On voit que \texttt{level: CRITICAL} déclenche automatiquement un tag \texttt{critical} dans Logstash. Le champ \texttt{failed\_attempts: 10} documente qu'on est bien au-delà du seuil de 5. Les champs \texttt{event\_type} et \texttt{alert\_type} sont les clés de voûte du pipeline : Logstash peut les lire directement sans passer par des expressions régulières complexes. C'est là toute l'importance du format JSON structuré — un log en texte libre forcerait à écrire des patterns Grok fragiles et coûteux à maintenir.

\subsubsection{Logstash - Enrichissement et tagging de sécurité}

Logstash est le cœur du pipeline. Son rôle est triple : parser le JSON reçu de Filebeat, enrichir chaque événement avec des tags de sécurité selon son type, puis indexer dans Elasticsearch avec le pattern \texttt{security-logs-YYYY.MM.dd} (rotation quotidienne des index).

\begin{table}[H]
  \centering
  \caption{Règles de tagging de sécurité appliquées par Logstash}
  \begin{tabularx}{\textwidth}{@{}llX@{}}
    \toprule
    \textbf{Condition} & \textbf{Tags ajoutés} & \textbf{alert\_type} \\
    \midrule
    \texttt{event\_type = brute\_force}          & brute\_force, security\_event, critical  & BRUTE\_FORCE \\
    \texttt{event\_type = auth\_failure}          & auth\_failure, security\_event           & AUTH\_FAILURE \\
    \texttt{event\_type = forbidden\_access}      & forbidden\_access, security\_event       & FORBIDDEN \\
    \texttt{event\_type = unauthorized\_access}   & unauthorized, security\_event            & UNAUTHORIZED \\
    \texttt{event\_type = suspicious\_upload}     & suspicious\_upload, security\_event      & SUSPICIOUS\_UPLOAD \\
    \texttt{event\_type = sqli\_attempt}          & sqli, security\_event, critical          & SQLI\_ATTEMPT \\
    \texttt{event\_type = suspicious\_ua}         & scanner\_ua, security\_event             & SUSPICIOUS\_UA \\
    \texttt{event\_type = geo\_anomaly}           & geo\_anomaly, security\_event            & GEO\_ANOMALY \\
    \texttt{status = 429}                         & rate\_limited, security\_event           & RATE\_LIMITED \\
    \texttt{level = ERROR / status = 500}         & server\_error                            & SERVER\_ERROR \\
    \texttt{service = api-gateway}                & gateway                                  & * \\
    \bottomrule
  \end{tabularx}
\end{table}

\medskip
Ces règles sont simples mais très efficaces. Chaque événement qui passe par Logstash repart avec des tags supplémentaires et, souvent, un champ \texttt{alert\_type} normalisé. Le tag \texttt{security\_event} est le plus important : il est systématiquement ajouté à tout événement suspect, ce qui permet au Dashboard de requêter \textit{tous} les incidents en une seule requête Elasticsearch, quelle que soit leur nature spécifique.

\subsubsection{Kibana - Exploration et investigation}

\screenshot{kibana_discover}{Kibana " interface Discover avec l'index security-logs-" et filtre event\_type: brute\_force}{kibana}

\medskip
L'interface Discover de Kibana offre une vue directe sur les données indexées. Ici, le filtre \texttt{event\_type: brute\_force} ne laisse passer que les tentatives d'attaque — chacune avec son IP source, son horodatage, et le nombre de tentatives échouées. Au-delà de la recherche ad hoc, Kibana permet de créer des visualisations (histogrammes, camemberts) et de les assembler en dashboards d'investigation, utiles lors d'une analyse post-incident.

\subsection{Monitoring Prometheus et Alertes}

\subsubsection{Architecture de scraping}

Prometheus scrape chaque cible toutes les \textbf{15 secondes} et évalue les règles d'alertes toutes les \textbf{30 secondes}. Cinq jobs de scraping sont configurés :

\begin{table}[H]
  \centering
  \caption{Cibles de scraping Prometheus (targets) et métriques clés collectées}
  \begin{tabularx}{\textwidth}{@{}llX@{}}
    \toprule
    \textbf{Job} & \textbf{Cible} & \textbf{Métriques clés} \\
    \midrule
    auth-service   & auth-service:5001/metrics  & \texttt{auth\_failures\_total}, \texttt{brute\_force\_total}, \texttt{flask\_http\_request\_total} \\
    api-service    & api-service:5002/metrics   & \texttt{forbidden\_access\_total}, \texttt{flask\_http\_request\_total} \\
    nginx-gateway  & nginx-exporter:9113        & \texttt{nginx\_http\_requests\_total}, \texttt{nginx\_connections\_active} \\
    node-exporter  & node-exporter:9100         & \texttt{node\_cpu\_seconds\_total}, \texttt{node\_memory\_*} \\
    alertmanager   & alertmanager:9093          & \texttt{alertmanager\_alerts}, \texttt{alertmanager\_notifications\_total} \\
    \bottomrule
  \end{tabularx}
\end{table}

\medskip
Chaque cible publie ses métriques au format texte Prometheus sur son endpoint \texttt{/metrics}. Prometheus les scrape régulièrement et les stocke dans sa base de données de séries temporelles (TSDB). C'est ce stockage qui rend PromQL si puissant : on peut calculer des taux sur une fenêtre glissante (\texttt{rate(...[5m])}), détecter des augmentations soudaines (\texttt{increase(...)}), et combiner ces indicateurs pour des alertes précises.

\screenshot{prometheus_targets}{Interface Prometheus - état des cibles de scraping, toutes UP}{prometheus_targets}

\medskip
La page Targets de Prometheus montre l'état de santé de chaque service surveillé. Tout est en vert — ce qui signifie que chaque endpoint \texttt{/metrics} répond correctement et dans les délais. La colonne \textit{Last Scrape} confirme que la collecte est bien continue. Si l'un de ces points passe au rouge (DOWN), la règle \texttt{ServiceDown} se déclenche dans la minute.

\subsubsection{Règles d'alertes}

8 règles d'alertes sont définies dans \texttt{monitoring/alert\_rules.yml}, réparties en deux groupes fonctionnels.

\textbf{Groupe Sécurité :}

\begin{table}[H]
  \centering
  \caption{Règles d'alertes du groupe Sécurité - expressions PromQL et seuils de déclenchement}
  \small
  \begin{tabularx}{\textwidth}{@{}p{3.5cm}>{\ttfamily\raggedright\arraybackslash}Xp{1.5cm}p{1.5cm}@{}}
    \toprule
    \textbf{Alerte} & \normalfont\textbf{Expression PromQL} & \textbf{Sév.} & \textbf{Délai} \\
    \midrule
    BruteForceDetected   & increase(brute\_force\_total[5m]) > 0                                                    & critical & immédiat \\
    HighAuthFailureRate  & rate(auth\_failures\_total[5m]) * 60 > 3                                                 & warning  & 2 min \\
    RateLimitViolations  & increase(flask\_http\_request\_total\newline\{status="429"\}[5m]) > 5                    & warning  & immédiat \\
    ForbiddenAccessSpike & rate(flask\_http\_request\_total\newline\{status="403"\}[5m]) * 60 > 5                  & warning  & 2 min \\
    \bottomrule
  \end{tabularx}
\end{table}

\medskip
Chacune de ces quatre règles cible un scénario d'attaque précis. \texttt{BruteForceDetected} est la plus urgente : elle se déclenche dès le premier événement brute force détecté (\texttt{increase > 0}), sans délai de grâce. \texttt{HighAuthFailureRate} est plus subtile — elle détecte une attaque diffuse à bas débit, où l'attaquant essaie lentement pour rester sous les radars du rate limiting. \texttt{ForbiddenAccessSpike} signale l'énumération de ressources protégées, souvent le signe d'une phase de reconnaissance préalable à une exploitation ciblée.

\textbf{Groupe Infrastructure :}

\begin{table}[H]
  \centering
  \caption{Règles d'alertes du groupe Infrastructure - expressions PromQL et seuils}
  \small
  \begin{tabularx}{\textwidth}{@{}p{3.5cm}>{\ttfamily\raggedright\arraybackslash}Xp{1.5cm}p{1.5cm}@{}}
    \toprule
    \textbf{Alerte} & \normalfont\textbf{Expression PromQL} & \textbf{Sév.} & \textbf{Délai} \\
    \midrule
    ServiceDown      & up == 0                                                                                          & critical & 1 min \\
    HighCPUUsage     & (1 - avg(rate(node\_cpu\_seconds\_total\newline\{mode="idle"\}[5m]))) * 100 > 80               & warning  & 2 min \\
    HighMemoryUsage  & (1 - node\_memory\_MemAvailable\_bytes\newline / node\_memory\_MemTotal\_bytes) * 100 > 85     & warning  & 2 min \\
    HighErrorRate    & rate(flask\_http\_request\_total\newline\{status=\textasciitilde"5.."\}[5m]) > 0.1             & warning  & 2 min \\
    \bottomrule
  \end{tabularx}
\end{table}

\medskip
Les règles infrastructure veillent sur la stabilité du système. La règle \texttt{ServiceDown} attend une minute avant de déclencher — un conteneur qui redémarre brièvement ne doit pas générer une alerte inutile. Les seuils CPU (80\%) et mémoire (85\%) sont calibrés pour alerter \textit{avant} la saturation, laissant du temps pour intervenir. Le calcul du CPU via le mode \texttt{idle} en complément à 1 est une technique Prometheus classique : elle donne le taux d'utilisation réel, indépendamment du nombre de cœurs.

\screenshot{prometheus_alerts}{Interface Prometheus - règle BruteForceDetected en état FIRING après simulation}{prometheus_alerts}

\medskip
Après la simulation brute force, la règle \texttt{BruteForceDetected} passe immédiatement en rouge avec le statut \textbf{FIRING}. Elle n'a pas de délai de grâce (\texttt{for: 0s}), donc la notification part vers Alertmanager dès le premier scraping qui détecte le compteur non nul. Les autres règles, si elles ont un délai configuré (\texttt{for: 2m}), passent d'abord par l'état PENDING avant de basculer en FIRING.

\subsubsection{Alertmanager - Routage et déduplication}

Alertmanager fait bien plus que simplement transmettre les alertes de Prometheus — il les traite intelligemment selon trois mécanismes :

\begin{itemize}[leftmargin=*]
  \item \textbf{Groupement} : plusieurs alertes de même nature sont fusionnées en une seule notification, évitant d'inonder l'équipe de messages redondants lors d'un incident.
  \item \textbf{Inhibition} : si un service est complètement down (\texttt{ServiceDown}), les alertes secondaires qui en découlent (CPU élevé, latence...) sont automatiquement supprimées — la cause racine est déjà connue.
  \item \textbf{Routage différencié} : les alertes critiques partent en 10 secondes, les alertes de sécurité en 15 secondes, avec des cadences de répétition différentes selon l'urgence.
\end{itemize}

\screenshot{alertmanager_ui}{Interface Alertmanager - alerte BruteForceDetected active avec ses labels et annotations}{alertmanager}

\medskip
L'interface affiche la carte d'alerte active avec ses labels (\texttt{severity: critical}, \texttt{service: auth-service}) et son annotation explicative. En production, on configurerait des receivers email, Slack ou PagerDuty pour que l'équipe de sécurité soit notifiée même sans surveiller activement le Dashboard.

\subsection{Tableaux de Bord Grafana}

Un avantage important de Grafana dans ce projet : les dashboards sont \textbf{auto-provisionnés} via des fichiers JSON. Pas de configuration manuelle, pas de clic dans l'interface — tout est prêt dès le \texttt{docker-compose up}.

\textbf{Dashboard Infrastructure :}

\screenshot{grafana_infra}{Grafana - Dashboard Infrastructure SecureWatch avec métriques système en temps réel}{grafana_infra}

\medskip
Ce premier dashboard surveille la santé de la machine hôte et des services : CPU, RAM, requêtes par seconde, temps de réponse au 95\textsuperscript{e} percentile, taux d'erreurs. La fenêtre temporelle de 30 minutes permet de repérer visuellement les pics d'activité et de les relier à ce qui se passait dans les logs au même moment.

\textbf{Dashboard Sécurité :}

\screenshot{grafana_security}{Grafana - Dashboard Sécurité SecureWatch avec courbes auth failures, 403 et 429}{grafana_security}

\medskip
Le second dashboard raconte une histoire différente : il affiche les signaux de sécurité. Compteur de brute force, courbes des 403, des 429 et des échecs d'authentification — tout est synchronisé sur la même fenêtre temporelle. Cette vue combinée est particulièrement révélatrice : quand les 403 et les auth failures montent simultanément, c'est souvent le signe qu'un attaquant explore l'architecture en parallèle de ses tentatives de login.

\subsection{Dashboard SecureWatch}

\subsubsection{Vue d'ensemble}

Le dashboard SecureWatch est une interface web développée en HTML/CSS/JavaScript pur, servie par un conteneur Nginx dédié. L'accès est protégé par une page de connexion : l'utilisateur s'authentifie via auth-service et reçoit un JWT contenant son rôle. Le token est stocké en \texttt{localStorage} et joint à chaque requête API. Le dashboard interroge Elasticsearch via un proxy nginx (évitant les problèmes CORS) et Prometheus via un second proxy, avec un rafraîchissement automatique toutes les \textbf{8 secondes}.

\screenshot{dashboard_overview}{Dashboard SecureWatch - vue d'ensemble avec compteurs de sécurité en temps réel}{dashboard}

\medskip
La page d'accueil du Dashboard offre un coup d'œil immédiat sur l'état de la sécurité : cinq compteurs en temps réel (total d'événements, échecs d'auth, accès interdits, erreurs système, brute force détectés), un graphique linéaire pour repérer les pics d'activité, et un donut qui visualise la répartition par type d'événement. Tout se rafraîchit toutes les 8 secondes automatiquement.

La barre supérieure intègre plusieurs fonctionnalités UX importantes : un \textbf{sélecteur de plage temporelle} (1h / 6h / 24h / 7j) qui ajuste dynamiquement toutes les requêtes Elasticsearch, ainsi qu'une \textbf{bascule de thème} (clair/sombre) dont la préférence est sauvegardée en \texttt{localStorage}. Le nom d'utilisateur connecté et son rôle sont affichés en temps réel, et un bouton de déconnexion ferme proprement la connexion SSE avant de nettoyer les tokens.

\begin{table}[H]
  \centering
  \caption{Permissions d'accès aux sections du Dashboard SecureWatch par rôle (RBAC)}
  \begin{tabular}{@{}lcccc@{}}
    \toprule
    \textbf{Section} & \textbf{Admin} & \textbf{Operator} & \textbf{User} \\
    \midrule
    Dashboard (vue d'ensemble) & \cmark & \cmark & \cmark \\
    Logs                       & \cmark & \cmark & \cmark \\
    Alertes                    & \cmark & \cmark & \xmark \\
    Carte géographique         & \cmark & \cmark & \xmark \\
    Services (état)            & \cmark & \cmark & \xmark \\
    Monitoring (infra + sécu)  & \cmark & \xmark & \xmark \\
    \bottomrule
  \end{tabular}
\end{table}

\medskip
Les restrictions de navigation sont appliquées côté client : les entrées de menu non autorisées sont grisées avec un attribut \texttt{aria-disabled}. Si un utilisateur tente d'y accéder directement, un toast d'accès refusé apparaît et la navigation est bloquée. Les restrictions côté serveur (RBAC sur les endpoints API) constituent la couche de protection primaire — la couche frontend n'est qu'un indicateur visuel.

\screenshot[0.90\textwidth]{ui_review/rbac_admin_dashboard}{Dashboard en mode Admin — toutes les sections accessibles, rôle affiché dans la barre supérieure}{rbac_admin}

\screenshot[0.90\textwidth]{ui_review/rbac_user1_dashboard}{Dashboard en mode User — sections Alertes, Services et Monitoring grisées dans la navigation}{rbac_user}

\subsubsection{Flux de Logs}

\screenshot{dashboard_logs}{Dashboard SecureWatch - flux de logs filtré sur le niveau CRITICAL}{logs}

\medskip
La section Logs présente les 1\,000 derniers événements significatifs — health checks et démarrages de services exclus, pour ne pas noyer les vraies informations. Un code couleur par niveau (INFO en bleu, WARNING en orange, ERROR en rouge, CRITICAL en rouge foncé) permet de scanner rapidement la liste. Les filtres par niveau et par service, combinés à la barre de recherche libre, transforment cette section en outil d'investigation ciblée. Un bouton \textbf{Export CSV} permet d'exporter les lignes filtrées actuellement visibles au format RFC-4180 avec un nom de fichier horodaté, utile pour des rapports d'incidents ou l'analyse forensique hors ligne.

\subsubsection{Alertes de Sécurité}

\screenshot{dashboard_alertes}{Dashboard SecureWatch - section Alertes après simulation brute force et forbidden scan}{alertes}

\medskip
La section Alertes agrège et présente les incidents de manière synthétique : une carte par type d'incident, colorée selon la criticité, avec l'IP source, le nombre d'occurrences et l'horodatage. La logique de détection est exécutée côté client, en interrogeant Elasticsearch à chaque rafraîchissement :

\begin{table}[H]
  \centering
  \caption{Règles de détection des alertes dans le Dashboard SecureWatch}
  \begin{tabularx}{\textwidth}{@{}llX@{}}
    \toprule
    \textbf{Type} & \textbf{Sévérité} & \textbf{Condition de déclenchement} \\
    \midrule
    Brute Force     & CRITIQUE & Présence d'un événement \texttt{brute\_force} pour une IP dans les logs \\
    Brute Force     & CRITIQUE & $\geq$ 5 échecs auth depuis la même IP en 24h \\
    Accès Interdit  & ÉLEVÉ    & $\geq$ 5 accès \texttt{forbidden\_access} depuis la même IP \\
    Erreurs Système & MOYEN    & $\geq$ 5 erreurs 500 sur le même service \\
    \bottomrule
  \end{tabularx}
\end{table}

\medskip
Ces règles sont intentionnellement simples pour être robustes. La double condition sur le brute force est un filet de sécurité : si pour une raison quelconque auth-service n'a pas encore généré son log CRITICAL, le comptage des échecs individuels prend le relais. Le seuil de 5 pour les accès interdits évite les faux positifs dus à des erreurs de navigation normales, tout en restant suffisamment bas pour détecter un scan rapide.

Chaque alerte peut être \textbf{ignorée individuellement} via un bouton « Ignorer » : une empreinte (\textit{fingerprint}) basée sur le type et l'IP source est stockée en \texttt{localStorage}. L'alerte n'apparaît plus lors des rafraîchissements suivants. Un bouton « Réinitialiser les alertes ignorées » permet de les restaurer. Cette persistance côté client évite les faux positifs visuels sans modifier les données Elasticsearch.

\screenshot[0.90\textwidth]{dashboard_review_4_geomap}{Dashboard SecureWatch — carte de géolocalisation des menaces (section Carte géographique)}{geomap}

\medskip
La section Carte géographique visualise les IPs sources détectées comme anomalies géographiques sur une carte mondiale SVG interactive. Chaque point représente une origine d'attaque simulée (nœuds Tor, plages IP à risque). Cette vue, accessible aux rôles admin et operator, offre une perspective spatiale des incidents complémentaire à la vue chronologique des logs.

\begin{infobox}[Note technique * Source des données 429]
Les réponses HTTP 429 sont générées par la gateway Nginx \textbf{avant} que la requête n'atteigne Flask. Par conséquent, \texttt{flask\_http\_request\_total\{status="429"\}} est toujours vide dans Prometheus. Les panneaux 429 du dashboard utilisent Elasticsearch (logs d'accès JSON de la gateway) comme source de données alternative.
\end{infobox}

\subsubsection{Monitoring et État des Services}

\screenshot{dashboard_monitoring_infra}{Dashboard SecureWatch onglet Monitoring, vue Infrastructure alimentée par Prometheus}{monitoring_infra}

\medskip
Cet onglet visualise les métriques système collectées par Prometheus depuis le \texttt{node-exporter} : CPU (jauge circulaire et courbe temporelle), RAM (barre de progression et courbe). Important à noter : ces métriques reflètent l'état de la \textit{machine hôte} physique, pas d'un conteneur en particulier. Pendant une simulation intensive, on peut donc voir le CPU grimper — et comprendre immédiatement d'où vient la charge.

\screenshot{dashboard_monitoring_security}{Dashboard SecureWatch onglet Monitoring, vue Sécurité}{monitoring_sec}

\medskip
L'onglet Sécurité creuse plus profond : compteurs absolus (brute force, 403, 429, 500) et trois graphiques temporels. Le plus intéressant est le graphique combiné qui superpose les 403 et les 429. Une corrélation forte entre ces deux courbes est un signal classique : un attaquant que le rate limiting ralentit tout en continuant à scanner des ressources protégées. Ce type de pattern serait très difficile à repérer sans visualisation temporelle.

\screenshot{dashboard_services}{Dashboard SecureWatch grille de statut en temps réel des 11 services}{services}

\medskip
La grille d'état sonde les 12 services (Redis inclus désormais) toutes les 8 secondes via des proxys nginx configurés dans le frontend (contournant ainsi les restrictions CORS du navigateur). Un point vert : le service répond à son endpoint de santé. Un point rouge : timeout ou erreur. C'est la vue la plus rapide pour détecter une panne — plus immédiat que de naviguer dans Prometheus ou de fouiller les logs Docker.

\screenshot[0.90\textwidth]{ui_review/theme_dark_dashboard}{Dashboard SecureWatch en thème sombre — vue d'ensemble}{theme_dark}

\screenshot[0.90\textwidth]{ui_review/theme_light_dashboard}{Dashboard SecureWatch en thème clair — vue d'ensemble}{theme_light}

\medskip
Le dashboard propose deux thèmes visuels (clair et sombre) sélectionnables depuis la barre de navigation. La préférence est sauvegardée en \texttt{localStorage} et restaurée automatiquement à chaque ouverture. Les deux thèmes couvrent l'intégralité des sections et composants, y compris les graphiques Chart.js dont les couleurs s'adaptent dynamiquement.

% ============================================================
%  CHAPITRE 4 * RÉSULTATS, TESTS ET ANALYSE
% ============================================================
\chapter{Résultats, Tests et Analyse}

\section{Introduction}

La théorie, c'est bien. La preuve, c'est mieux. Ce chapitre met SecureWatch à l'épreuve : six scénarios d'attaque simulés, des résultats mesurés, et une analyse honnête de ce qui fonctionne bien — et de ce qui peut encore être amélioré.

\section{Démonstration des Cas d'Usage}

\subsection{Script de simulation d'attaques}

Le fichier \texttt{scripts/generate\_attacks.py} implémente 6 scénarios d'attaque ciblant la gateway via HTTPS (\texttt{https://localhost:8443}). Chaque scénario peut être lancé individuellement ou en séquence complète :

\begin{lstlisting}[style=bash, caption={Lancement de l'ensemble des scénarios d'attaque}]
python scripts/generate_attacks.py --scenario all
\end{lstlisting}

\medskip
Une seule commande déclenche les 6 scénarios en séquence. On peut aussi les lancer individuellement en remplaçant \texttt{all} par le nom du scénario ciblé (\texttt{brute-force}, \texttt{forbidden-scan}, etc.). Une fois le script terminé, il faut patienter 30 à 60 secondes avant de consulter les interfaces — le temps que Filebeat transmette les logs à Logstash et que Prometheus ait évalué ses règles. Se précipiter sur Kibana 2 secondes après l'exécution ne donnera rien.

\begin{infobox}[Note HTTPS - Certificat auto-signé]
Depuis l'implémentation de TLS sur la gateway, le script utilise \texttt{verify=False} pour accepter le certificat auto-signé (comportement attendu en développement), et les alertes \texttt{InsecureRequestWarning} de la bibliothèque \texttt{urllib3} sont supprimées pour un affichage propre dans le terminal.
\end{infobox}

\subsection{Scénario 1 - Attaque Brute Force}

\textbf{Objectif :} simuler une attaque par dictionnaire sur \texttt{/auth/login} depuis plusieurs adresses IP d'attaquants.

\begin{lstlisting}[style=bash, caption={Lancement du scénario brute force}]
python scripts/generate_attacks.py --scenario brute-force
\end{lstlisting}

\medskip
Ce scénario simule une attaque par dictionnaire classique : des requêtes POST répétées vers \texttt{/auth/login} avec des credentials incorrects, depuis plusieurs IPs simulées. Les tentatives sont envoyées en rafale rapide, sans délai entre elles — exactement comme un vrai outil de force brute. Dès que le seuil de 5 échecs est atteint depuis la même IP, auth-service génère un log CRITICAL et incrémente le compteur \texttt{brute\_force\_total} dans Prometheus.

\textbf{Déroulement :}
\begin{itemize}[leftmargin=*]
  \item 8 à 12 tentatives avec credentials invalides depuis des IPs simulées (ex: \texttt{185.220.101.5}, \texttt{10.0.0.5}, \texttt{192.168.1.42})
  \item Chaque tentative : HTTP 401 (Unauthorized) + log \texttt{auth\_failure} dans ELK
  \item Après $\geq$ 5 échecs depuis la même IP : log \texttt{brute\_force} niveau CRITICAL + incrémentation de \texttt{brute\_force\_total} dans Prometheus
\end{itemize}

\textbf{Résultats observables :}
\begin{itemize}[leftmargin=*]
  \item[\cmark] Alerte \textbf{BRUTE FORCE CRITIQUE} dans le Dashboard (section Alertes) avec IP source
  \item[\cmark] Règle \texttt{BruteForceDetected} en état FIRING dans Prometheus (délai $\leq$ 30 s)
  \item[\cmark] Alerte active dans Alertmanager avec labels \texttt{severity: critical}
  \item[\cmark] Logs CRITICAL visibles dans Kibana (\texttt{event\_type: brute\_force})
  \item[\cmark] Compteur \texttt{brute\_force\_total} incrémenté dans les métriques Prometheus
\end{itemize}

\subsection{Scénario 2 - Scan de Routes Interdites}

\textbf{Objectif :} simuler un attaquant cherchant des ressources sensibles par énumération de chemins.

\begin{lstlisting}[style=bash]
python scripts/generate_attacks.py --scenario forbidden-scan
\end{lstlisting}

\medskip
Avant d'attaquer, un attaquant explore. Ce scénario simule cette phase de reconnaissance : l'énumération de chemins pour découvrir des ressources d'administration oubliées ou mal protégées. C'est un des patterns les plus courants dans les logs de serveurs web en production. Un 403 en réponse révèle qu'une ressource existe mais est protégée — c'est une information précieuse pour l'attaquant. Un 404 indique un chemin inexistant.

Routes scannées : \texttt{/admin}, \texttt{/config}, \texttt{/secrets}, \texttt{/users}, \texttt{/debug}, \texttt{/env}, \texttt{/backup}, \texttt{/keys}, \texttt{/tokens}.

\textbf{Résultats observables :}
\begin{itemize}[leftmargin=*]
  \item[\cmark] HTTP 403 pour les routes protégées, HTTP 404 pour les routes inexistantes
  \item[\cmark] Logs \texttt{forbidden\_access} dans ELK avec chemin tenté et IP source
  \item[\cmark] Alerte \textbf{ACCÈS INTERDIT ÉLEVÉ} dans le Dashboard après $\geq$ 5 occurrences par IP
  \item[\cmark] Règle \texttt{ForbiddenAccessSpike} en état PENDING puis FIRING dans Prometheus
\end{itemize}

\subsection{Scénario 3 - Test Rate Limiting}

\textbf{Objectif :} démontrer le rate limiting de la gateway en saturant les zones \texttt{auth\_zone} et \texttt{api\_zone}.

\begin{lstlisting}[style=bash]
python scripts/generate_attacks.py --scenario rate-limit
\end{lstlisting}

\medskip
Ce test vérifie que les zones de rate limiting fonctionnent comme prévu. Il se déroule en deux phases pour confirmer l'indépendance des seuils : d'abord la zone API (plus permissive), puis la zone Auth (très restrictive). L'objectif est de voir les 429 apparaître au bon moment, accompagnés du corps JSON de la gateway — preuve que le blocage se fait bien au bon niveau.

\textbf{Phase 1 - Zone API :} 60 requêtes vers \texttt{/api/data} (limite 30/min, burst 50). Les HTTP 429 apparaissent à partir de la 38\textsuperscript{e} requête (burst absorbé, puis seau plein).

\textbf{Phase 2 - Zone Auth :} 10 tentatives de login invalides (limite 5/min, burst 10). Les HTTP 429 apparaissent dès la 2\textsuperscript{e} requête (burst très réduit sur la zone auth, plus restrictive).

\screenshot[0.85\textwidth]{rate_limit_terminal}{Sortie terminal du scénario rate-limit phase 1 (API) et phase 2 (Auth) avec requêtes bloquées}{ratelimit_term}

\medskip
La sortie terminal raconte exactement ce qui se passe : d'abord des 200 normaux pendant l'absorption du burst, puis des 429 en série une fois le seau plein. La phase Auth est nettement plus agressive dans le blocage — le burst plus petit laisse beaucoup moins de marge. Ce comportement différencié par zone montre que la configuration Nginx est correctement appliquée : les deux zones opèrent indépendamment, avec leurs propres seuils.

\subsection{Scénario 4 - Injection SQL et User-Agent offensif}

\textbf{Objectif :} démontrer la détection de tentatives d'injection SQL dans les champs d'authentification et l'identification d'outils d'attaque via leur User-Agent.

\begin{lstlisting}[style=bash]
python scripts/generate_attacks.py --scenario sqli
\end{lstlisting}

\medskip
Ce scénario envoie des requêtes de login avec des payloads d'injection SQL classiques (\texttt{' OR '1'='1}, \texttt{UNION SELECT}, opérateurs hexadécimaux) et des requêtes avec des User-Agents typiques d'outils d'attaque (sqlmap, nikto, gobuster). Chaque détection génère un log avec le payload tronqué pour éviter toute réinjection dans les systèmes de log.

\textbf{Résultats observables :}
\begin{itemize}[leftmargin=*]
  \item[\cmark] HTTP 400 retourné pour chaque tentative d'injection SQL (\texttt{event\_type: sqli\_attempt})
  \item[\cmark] Compteur \texttt{sqli\_attempts\_total} incrémenté dans Prometheus
  \item[\cmark] Logs \texttt{suspicious\_ua} générés pour les User-Agents offensifs
  \item[\cmark] Compteur \texttt{scanner\_ua\_total} incrémenté dans Prometheus
  \item[\cmark] Événements visibles dans Kibana et le flux de logs du Dashboard
\end{itemize}

\subsection{Scénario 5 - Anomalie géographique}

\textbf{Objectif :} simuler des requêtes provenant d'adresses IP associées à des origines géographiques à risque (nœuds Tor, plages d'États connus pour les cyberattaques).

\begin{lstlisting}[style=bash]
python scripts/generate_attacks.py --scenario geo-anomaly
\end{lstlisting}

\medskip
Le script injecte un header \texttt{X-Forwarded-For} avec des IPs simulées appartenant à des plages surveillées : nœuds de sortie Tor (\texttt{185.220.x.x}, \texttt{162.247.x.x}), plages russes, chinoises et iraniennes. Api-service analyse le premier hop et génère un log \texttt{geo\_anomaly} avec le code pays et le nom de la région identifiée.

\textbf{Résultats observables :}
\begin{itemize}[leftmargin=*]
  \item[\cmark] Logs \texttt{geo\_anomaly} dans ELK avec \texttt{country\_code} et \texttt{country\_name}
  \item[\cmark] Compteur \texttt{geo\_anomaly\_total} incrémenté dans Prometheus
  \item[\cmark] Points affichés sur la carte géographique du Dashboard (section dédiée)
  \item[\cmark] Événements filtrables dans Kibana par \texttt{event\_type: geo\_anomaly}
\end{itemize}

\screenshot[0.90\textwidth]{dashboard_review_4b_geomap_table}{Dashboard SecureWatch — liste des anomalies géographiques détectées avec pays et IPs sources}{geomap_table}

\subsection{Résultats Consolidés}

Après exécution de l'ensemble des scénarios de simulation, Elasticsearch a indexé plus de 4\,400 événements :

\begin{table}[H]
  \centering
  \caption{Distribution des événements indexés dans Elasticsearch après simulation complète des 6 scénarios}
  \begin{tabular}{@{}lr@{}}
    \toprule
    \textbf{Type d'événement} & \textbf{Occurrences} \\
    \midrule
    \texttt{health\_check}      & 2\,745 \\
    \texttt{gateway\_request}   & 888 \\
    \texttt{api\_access}        & 274 \\
    \texttt{auth\_failure}      & 147 \\
    \texttt{service\_start}     & 92 \\
    \texttt{brute\_force}       & 87 \\
    \texttt{not\_found}         & 80 \\
    \texttt{suspicious\_upload} & 55 \\
    \texttt{auth\_success}      & 39 \\
    \texttt{forbidden\_access}  & 26 \\
    \texttt{server\_error}      & 26 \\
    \texttt{unauthorized\_access} & 15 \\
    \midrule
    \textbf{Total}              & \textbf{4\,474+} \\
    \bottomrule
  \end{tabular}
\end{table}

\medskip
Ces chiffres appellent quelques explications. Les 2\,745 \texttt{health\_check} dominent le total — ce sont les sondages automatiques de Prometheus et du Dashboard, toutes les 8 à 15 secondes pendant toute la durée de la session. Ils sont systématiquement filtrés des vues de sécurité (via \texttt{must\_not} dans les requêtes Elasticsearch) pour ne pas noyer les vraies informations. Les 87 événements \texttt{brute\_force} et 147 \texttt{auth\_failure} sont les plus importants : ils confirment que les incidents critiques ont bien été générés, transmis et indexés sans perte.

\section{Analyse de Sécurité}

\subsection{Résilience face aux accès non autorisés}

L'évaluation de la résilience de SecureWatch face aux différents vecteurs d'attaque peut être résumée comme suit :

\begin{table}[H]
  \centering
  \small
  \setlength{\tabcolsep}{5pt}
  \caption{Évaluation de la résilience de SecureWatch par type d'attaque}
  \begin{tabularx}{\textwidth}{@{}p{3.0cm}Xp{2.0cm}p{2.4cm}@{}}
    \toprule
    \textbf{Vecteur} & \textbf{Mécanisme de défense} & \textbf{Détection} & \textbf{Blocage} \\
    \midrule
    Brute force
      & Rate limiting \texttt{auth\_zone} + détection applicative + Redis
      & \cmark{} Immédiate & \cmark{} Partiel \\
    \addlinespace
    Scan ressources (403)
      & Contrôles RBAC + journalisation gateway
      & \cmark{} Seuil & \xmark{} Non bloqué \\
    \addlinespace
    Accès rôle insuffisant
      & RBAC serveur (JWT claims) + nav frontend restreinte
      & \cmark{} Immédiate & \cmark{} 401/403 \\
    \addlinespace
    Injection SQL
      & Regex multi-patterns \texttt{@before\_request} (auth)
      & \cmark{} Immédiate & \cmark{} HTTP 400 \\
    \addlinespace
    Outil offensif (UA)
      & Regex User-Agent : sqlmap, nikto, gobuster...
      & \cmark{} Immédiate & \xmark{} Log seul \\
    \addlinespace
    Origine géo suspecte
      & Préfixe IP : Tor, Russie, Chine, Iran...
      & \cmark{} Immédiate & \xmark{} Log seul \\
    \addlinespace
    Flood / DDoS
      & Rate limiting \texttt{global\_zone} (100 req/min)
      & \cmark{} Via 429 & \cmark{} Oui \\
    \addlinespace
    Accès direct services
      & Isolation réseau — aucun port exposé sur l'hôte
      & --- & \cmark{} Total \\
    \addlinespace
    Écoute réseau (MITM)
      & TLS 1.2/1.3 + HSTS
      & --- & \cmark{} Total \\
    \addlinespace
    Upload malveillant
      & RBAC (admin/operator) + vérification MIME
      & \cmark{} Immédiate & \cmark{} HTTP 415 \\
    \bottomrule
  \end{tabularx}
\end{table}

\medskip
Ce tableau est une lecture honnête des forces et faiblesses. Le point fort le plus solide reste l'isolation réseau : il n'existe tout simplement aucun chemin direct vers les microservices depuis l'extérieur. TLS protège toutes les communications. Le rate limiting est efficace contre les attaques à fort volume. La limite principale est réelle : un attaquant patient qui distribue ses tentatives sur plusieurs IPs en respectant les délais peut passer sous les radars. Une liste noire IP dynamique (intégration fail2ban, par exemple) serait nécessaire pour fermer cette brèche en production.

\subsection{Efficacité du pipeline de détection}

Le pipeline complet \textit{attaque -> log -> Elasticsearch -> alerte Dashboard} fonctionne en \textbf{moins de 30 secondes} dans les conditions normales de fonctionnement :

\begin{itemize}[leftmargin=*]
  \item \textbf{Génération du log} : instantanée (dans la même transaction HTTP)
  \item \textbf{Collecte Filebeat -> Logstash} : 1--3 secondes (intervalle de scraping Filebeat)
  \item \textbf{Indexation Elasticsearch} : moins de 1 seconde
  \item \textbf{Détection Dashboard} : jusqu'à 8 secondes (intervalle de rafraîchissement)
  \item \textbf{Évaluation règle Prometheus} : jusqu'à 30 secondes (intervalle d'évaluation)
\end{itemize}

Moins de 30 secondes de bout en bout, c'est tout à fait acceptable pour un système de détection en temps quasi-réel. Si on voulait aller plus vite, on pourrait réduire l'intervalle de collecte de Filebeat (500 ms plutôt que quelques secondes) et passer Prometheus à 10 secondes de scraping. Mais dans un contexte pédagogique, la configuration actuelle suffit largement.

\section{Limites et Perspectives}

\subsection{Difficultés rencontrées}

Aucun projet n'avance sans heurts. Voici quatre problèmes que nous avons rencontrés et résolus :

\begin{itemize}[leftmargin=*]
  \item \textbf{Graphiques Chart.js invisibles} : les canvas dans des sections masquées (\texttt{display:none}) ont des dimensions nulles à l'initialisation — Chart.js ne peut pas dessiner. Solution : initialisation paresseuse (\textit{lazy init}), déclenchée uniquement quand l'utilisateur navigue vers la section Monitoring.
  \item \textbf{Métriques 429 absentes de Prometheus} : Nginx bloque les requêtes rate-limitées avant qu'elles n'atteignent Flask — le compteur \texttt{flask\_http\_request\_total\{status="429"\}} reste à zéro. Solution : on interroge directement Elasticsearch (logs d'accès de la gateway) pour les données 429, contournant le problème à la source.
  \item \textbf{Cache DNS du frontend Nginx} : après recréation d'un conteneur, Nginx conserve l'ancienne IP résolue pour Prometheus, causant des échecs de proxy. Solution : un simple \texttt{nginx -s reload} dans le conteneur frontend force la résolution DNS.
  \item \textbf{Flood de health\_check dans ELK} : 2\,745 health checks occupaient les 1\,000 premiers résultats, masquant complètement les événements de sécurité. Solution : exclusion systématique de ces types via \texttt{must\_not} dans toutes les requêtes du Dashboard.
\end{itemize}

\subsection{Contraintes liées au déploiement local}

Faire tourner 13 services sur un poste de développement a ses contraintes :

\begin{itemize}[leftmargin=*]
  \item \textbf{RAM} : Elasticsearch seul consomme 2 Go de heap JVM. La stack complète demande 6 à 8 Go, ce qui impose une machine avec au moins 12 Go de RAM disponible — pas toujours évident sur un laptop étudiant.
  \item \textbf{Démarrage lent} : Elasticsearch et Logstash mettent 30 à 60 secondes avant d'être opérationnels. Sans \texttt{healthcheck} et \texttt{depends\_on} bien configurés dans Docker Compose, les services qui en dépendent démarrent trop tôt et échouent silencieusement.
  \item \textbf{Certificat auto-signé} : les navigateurs affichent un avertissement de sécurité, et les clients HTTP doivent désactiver la vérification TLS en développement. Acceptable localement, à corriger impérativement en production.
\end{itemize}

\subsection{Perspectives d'évolution}

\begin{itemize}[leftmargin=*]
  \item \textbf{Notifications Alertmanager} : configurer des receivers email, Slack ou Discord pour les alertes critiques, assurant une notification des équipes de sécurité en dehors de l'interface web.
  \item \textbf{Certificat TLS signé par CA} : remplacer le certificat auto-signé par Let's Encrypt (via cert-manager en Kubernetes) ou une CA interne d'entreprise.
  \item \textbf{Blocage IP dynamique} : intégrer fail2ban ou une liste noire IP Nginx mise à jour automatiquement en fonction des détections de brute force dans Elasticsearch.
  \item \textbf{Authentification dashboard} : ajouter une couche d'authentification (JWT ou session HTTP) sur le Dashboard SecureWatch lui-même, pour éviter que n'importe qui sur le réseau puisse accéder aux informations de sécurité.
  \item \textbf{Machine learning} : intégrer des modèles de détection d'anomalies sur les séries temporelles Prometheus (ex: LSTM pour les pics de trafic anormaux).
  \item \textbf{Déploiement cloud} : porter l'architecture sur Kubernetes avec Helm charts pour un déploiement scalable en production, avec secrets gérés par Vault ou AWS Secrets Manager.
  \item \textbf{Tests de charge automatisés} : intégrer Locust ou k6 pour calibrer et valider les zones de rate limiting sous charge réelle.
\end{itemize}

% ============================================================
%  CONCLUSION GÉNÉRALE
% ============================================================
\chapter*{Conclusion générale}
\addcontentsline{toc}{chapter}{Conclusion générale}
\markboth{Conclusion générale}{Conclusion générale}

\begin{successbox}[Bilan des réalisations]
\begin{itemize}[leftmargin=*]
  \item[\cmark] \textbf{Pipeline de logs complet} : Filebeat -> Logstash (règles de tagging sécurité) -> Elasticsearch -> Kibana, avec rotation quotidienne et rétention automatique ILM à 14 jours.
  \item[\cmark] \textbf{API Gateway HTTPS} : terminaison TLS, rate limiting sur 3 zones (auth, api, global), logs JSON structurés, isolation réseau totale des microservices.
  \item[\cmark] \textbf{RBAC 3 niveaux} : contrôle d'accès admin/operator/user appliqué côté serveur (JWT claims) et côté client (navigation conditionnelle + toasts d'accès refusé).
  \item[\cmark] \textbf{Alertes temps réel SSE} : endpoint \texttt{/events/stream} sur auth-service, notifications toast immédiates sur brute force sans attendre le polling.
  \item[\cmark] \textbf{Redis persistence} : compteurs brute-force et alertes webhook survivent aux redémarrages, avec fallback mémoire automatique.
  \item[\cmark] \textbf{Monitoring temps réel} : Prometheus (5 targets, 8 règles d'alertes), Grafana (2 dashboards auto-provisionnés), Alertmanager (webhook receiver actif vers api-service).
  \item[\cmark] \textbf{Dashboard interactif} : 6 sections, RBAC nav, sélecteur plage temporelle, export CSV, gestion des thèmes, carte géographique des menaces, détection d'incidents en temps réel.
  \item[\cmark] \textbf{Simulation d'attaques} : 8 scénarios couvrant brute force, forbidden scan, rate limiting, accès non autorisés, uploads suspects, injection SQL, User-Agents offensifs, anomalies géographiques.
  \item[\cmark] \textbf{Architecture reproductible} : 14 services orchestrés via Docker Compose sur 3 réseaux isolés, démarrage complet en une commande.
\end{itemize}
\end{successbox}

SecureWatch prouve une chose simple mais importante : on peut construire une plateforme de sécurité sérieuse avec zéro coût de licence. Elasticsearch, Prometheus, Grafana, Nginx — des outils open-source maintenus par des milliers de contributeurs, déployables en une commande, et capables de répondre à de vrais besoins de surveillance.

Ce projet a été, avant tout, une expérience d'apprentissage intense. Configurer TLS à la main, calibrer un rate limiting zone par zone, câbler un pipeline de logs de bout en bout, écrire des règles d'alertes PromQL — ce sont des compétences qu'on ne peut pas acquérir en lisant un manuel. La simulation d'attaques réelles a confirmé que les mécanismes fonctionnent. Elle a aussi révélé leurs limites, ce qui est tout aussi précieux : un système dont on ne connaît pas les faiblesses est dangereux. Il illustre concrètement le principe de \textit{défense en profondeur} : aucun mécanisme seul ne suffit, mais leur combinaison — réseau, protocole, applicatif, surveillance — offre une résilience réelle.

Le chemin qui mène de ce projet vers une plateforme de niveau production est tracé : Kubernetes pour l'orchestration, Let's Encrypt pour les certificats, fail2ban pour le blocage IP dynamique, Vault pour la gestion des secrets, et peut-être un jour de la détection d'anomalies par machine learning sur les séries temporelles Prometheus. La base architecturale est solide. Il ne reste qu'à continuer de construire dessus.

% ── Bibliographie ─────────────────────────────────────────────
\begin{thebibliography}{99}

\bibitem{elk}
Elastic N.V. (2024).
\textit{Elasticsearch Reference 8.12}.
\url{https://www.elastic.co/guide/en/elasticsearch/reference/8.12/index.html}

\bibitem{logstash}
Elastic N.V. (2024).
\textit{Logstash Reference 8.12 * Configuring Logstash}.
\url{https://www.elastic.co/guide/en/logstash/8.12/index.html}

\bibitem{filebeat}
Elastic N.V. (2024).
\textit{Filebeat Reference 8.12}.
\url{https://www.elastic.co/guide/en/beats/filebeat/8.12/index.html}

\bibitem{prometheus}
Prometheus Authors (2024).
\textit{Prometheus Documentation * Overview and PromQL}.
\url{https://prometheus.io/docs/introduction/overview/}

\bibitem{alertmanager}
Prometheus Authors (2024).
\textit{Alertmanager Documentation}.
\url{https://prometheus.io/docs/alerting/latest/alertmanager/}

\bibitem{nginx}
NGINX, Inc. (2024).
\textit{NGINX Documentation * ngx\_http\_limit\_req\_module (Rate Limiting)}.
\url{https://nginx.org/en/docs/http/ngx_http_limit_req_module.html}

\bibitem{docker}
Docker, Inc. (2024).
\textit{Docker Compose Documentation * Compose file reference}.
\url{https://docs.docker.com/compose/}

\bibitem{grafana}
Grafana Labs (2024).
\textit{Grafana Documentation * Dashboard Provisioning}.
\url{https://grafana.com/docs/grafana/latest/administration/provisioning/}

\bibitem{flask}
Pallets Projects (2024).
\textit{Flask Documentation}.
\url{https://flask.palletsprojects.com/}

\bibitem{jwt}
Jones, M. et al. (2015).
\textit{JSON Web Token (JWT) * RFC 7519}.
\url{https://datatracker.ietf.org/doc/html/rfc7519}

\bibitem{owasp}
OWASP Foundation (2024).
\textit{OWASP API Security Top 10}.
\url{https://owasp.org/www-project-api-security/}

\bibitem{verizon}
Verizon (2024).
\textit{Data Breach Investigations Report (DBIR) 2024}.
\url{https://www.verizon.com/business/resources/reports/dbir/}

\end{thebibliography}

% ── Annexes ───────────────────────────────────────────────────
\appendix

\chapter{Extrait du Fichier Docker Compose}

Le fichier \texttt{docker-compose.yml} définit l'ensemble des 13 services et leurs configurations. L'extrait ci-dessous présente la définition des services principaux :

\begin{lstlisting}[style=yaml, caption={Extrait docker-compose.yml * API Gateway et Auth Service}]
version: "3.8"

services:
  gateway:
    image: nginx:alpine
    ports:
      - "8080:8080"   # HTTP -> redirect HTTPS
      - "8443:8443"   # HTTPS
    volumes:
      - ./gateway/nginx.conf:/etc/nginx/nginx.conf:ro
      - ./certs:/etc/nginx/certs:ro
      - logs_volume:/logs
    networks:
      - gateway-net
      - monitoring-net
      - elk
    depends_on:
      - auth-service
      - api-service

  auth-service:
    build: ./auth-service
    expose:
      - "5001"         # Port interne uniquement, pas expose sur l'hote
    volumes:
      - logs_volume:/logs
    networks:
      - gateway-net
      - monitoring-net
      - elk

  elasticsearch:
    image: docker.elastic.co/elasticsearch/elasticsearch:8.12.0
    environment:
      - discovery.type=single-node
      - xpack.security.enabled=false
      - "ES_JAVA_OPTS=-Xms1g -Xmx1g"
    networks:
      - elk
    healthcheck:
      test: ["CMD", "curl", "-f", "http://localhost:9200"]
      interval: 30s
      timeout: 10s
      retries: 5
\end{lstlisting}

\medskip
Cet extrait révèle les choix de sécurité encodés dans la configuration Docker Compose. Le point le plus important : seule la \texttt{gateway} utilise \texttt{ports} (mapping hôte:conteneur). Les microservices utilisent \texttt{expose}, qui rend le port accessible uniquement aux autres conteneurs du même réseau — jamais à l'hôte. \texttt{depends\_on} garantit l'ordre de démarrage, mais pas la disponibilité applicative — d'où le \texttt{healthcheck} d'Elasticsearch : Logstash et Kibana attendent que le moteur réponde sur son port HTTP avant de tenter de s'y connecter. La contrainte \texttt{ES\_JAVA\_OPTS=-Xms1g -Xmx1g} fixe le heap Java à exactement 1 Go, évitant les réallocations JVM et la saturation mémoire sur une machine de développement.

\chapter{Configuration Logstash - Pipeline de Sécurité}

Le fichier \texttt{logstash/pipeline/logstash.conf} définit le pipeline de traitement des logs :

\begin{lstlisting}[style=yaml, caption={Configuration Logstash * input, filter et règles de tagging sécurité}]
input {
  beats {
    port => 5044
  }
}

filter {
  json {
    source => "message"
  }

  # Normalisation du timestamp
  date {
    match => ["timestamp", "ISO8601"]
    target => "@timestamp"
  }

  # Tagging securite par type d'evenement
  if [event_type] == "brute_force" {
    mutate {
      add_tag => ["brute_force", "security_event", "critical"]
      add_field => { "alert_type" => "BRUTE_FORCE" }
    }
  } else if [event_type] == "auth_failure" {
    mutate {
      add_tag => ["auth_failure", "security_event"]
      add_field => { "alert_type" => "AUTH_FAILURE" }
    }
  } else if [event_type] == "forbidden_access" {
    mutate {
      add_tag => ["forbidden_access", "security_event"]
      add_field => { "alert_type" => "FORBIDDEN" }
    }
  } else if [status] == "429" {
    mutate {
      add_tag => ["rate_limited", "security_event"]
      add_field => { "alert_type" => "RATE_LIMITED" }
    }
  }
}

output {
  elasticsearch {
    hosts => ["elasticsearch:9200"]
    index => "security-logs-%{+YYYY.MM.dd}"
  }
}
\end{lstlisting}

\medskip
La structure en trois blocs de Logstash est simple à lire. Le bloc \texttt{input} écoute sur le port TCP 5044 et attend les données envoyées par Filebeat via le protocole Beats. Le bloc \texttt{filter} applique deux transformations successives : le plugin \texttt{json} déstructure le message brut en champs Elasticsearch distincts, puis le plugin \texttt{date} normalise le timestamp applicatif au format \texttt{@timestamp} standard, ce qui garantit un tri chronologique correct dans Kibana. La chaîne \texttt{if/else if} est le moteur de classification — chaque type d'événement reçoit ses tags et son \texttt{alert\_type}. Le bloc \texttt{output} indexe dans Elasticsearch avec une rotation quotidienne (\texttt{YYYY.MM.dd}) : un nouvel index est créé automatiquement chaque jour, ce qui simplifie la rétention et distribue la charge d'indexation.

\chapter{Configuration Nginx - API Gateway}

Extrait de la configuration Nginx pour la gateway sécurisée avec rate limiting et logs JSON :

\begin{lstlisting}[style=yaml, caption={Extrait nginx.conf * zones de rate limiting et en-têtes de sécurité}]
http {
  # Zones de rate limiting (algorithme leaky bucket)
  limit_req_zone $binary_remote_addr
    zone=auth_zone:10m rate=5r/m;
  limit_req_zone $binary_remote_addr
    zone=api_zone:10m  rate=30r/m;
  limit_req_zone $binary_remote_addr
    zone=global_zone:10m rate=100r/m;

  # Format de log JSON structure
  log_format json_log escape=json
    '{"timestamp":"$time_iso8601",'
    '"service":"api-gateway",'
    '"ip":"$remote_addr",'
    '"method":"$request_method",'
    '"uri":"$request_uri",'
    '"status":$status,'
    '"request_time":"$request_time",'
    '"event_type":"gateway_request"}';

  server {
    listen 8443 ssl;
    ssl_certificate     /etc/nginx/certs/server.crt;
    ssl_certificate_key /etc/nginx/certs/server.key;

    # En-tetes de securite
    add_header Strict-Transport-Security
      "max-age=31536000; includeSubDomains" always;
    add_header X-Frame-Options "DENY" always;
    add_header X-Content-Type-Options "nosniff" always;

    # Route authentification avec rate limiting strict
    location /auth/ {
      limit_req zone=auth_zone burst=10 nodelay;
      limit_req_status 429;
      proxy_pass http://auth-service:5001/;
    }

    # Route API avec rate limiting standard
    location /api/ {
      limit_req zone=api_zone burst=50 nodelay;
      limit_req_status 429;
      proxy_pass http://api-service:5002/;
    }
  }
}
\end{lstlisting}

\medskip
Trois couches de sécurité transparaissent dans cette configuration. La première : le \textbf{rate limiting}. La directive \texttt{limit\_req\_zone} définit chaque zone avec sa clé de partitionnement (\texttt{\$binary\_remote\_addr}, l'IP en binaire — 4 fois plus compact qu'une chaîne), sa taille mémoire (10 Mo, suffisant pour environ 160\,000 IPs distinctes), et son débit. La deuxième : le \textbf{format de log JSON}. Le paramètre \texttt{escape=json} garantit que les caractères spéciaux dans les URLs sont correctement échappés — sans ça, un seul guillemet dans une URL casserait le JSON et ferait échouer Logstash. La troisième : la \textbf{configuration TLS} avec ses en-têtes de sécurité. Sur les blocs \texttt{location}, l'option \texttt{nodelay} est importante : sans elle, Nginx met les requêtes du burst en file d'attente, introduisant une latence artificielle. Avec \texttt{nodelay}, elles passent immédiatement ou reçoivent un 429 instantané — le comportement attendu pour une gateway de sécurité.

\end{document}